\documentclass[11pt,a4paper]{article}

% ── Packages ──────────────────────────────────────────────────────────────
\usepackage[utf8]{inputenc}
\usepackage[T1]{fontenc}
\usepackage{amsmath,amssymb,amsthm}
\usepackage{amscd}
\usepackage{booktabs}
\usepackage{graphicx}
\usepackage{hyperref}
\usepackage[margin=1in]{geometry}
\usepackage{enumitem}
\usepackage{xcolor}
\definecolor{hlblue}{RGB}{0,80,160}
\hypersetup{colorlinks=true,linkcolor=hlblue,urlcolor=hlblue,citecolor=hlblue}

% ── Theorems ──────────────────────────────────────────────────────────────
\newtheorem{principle}{Principle}[section]
\newtheorem{theorem}{Theorem}[section]
\newtheorem{definition}{Definition}[section]
\newtheorem{proposition}{Proposition}[section]
\theoremstyle{remark}
\newtheorem{remark}{Remark}[section]

% ── Title ─────────────────────────────────────────────────────────────────
\title{\textbf{Emerging World Models}\\[4pt]
\large A Structural Framework for Self-Organizing World Representations
\\[4pt]
\large \textsc{Consolidation Document --- August 2026}}
\date{}

\begin{document}
\maketitle

\begin{abstract}
We propose \textbf{Emerging World Models} (EWM)---a structural framework for world
representations whose ontology is not predefined but crystallizes through measurement.
EWM is governed by four foundation principles. \textbf{(1) Emergent Ontology:} elements
and subsystems are not primitive; they emerge from the accumulation of IICA measurement
interactions. \textbf{(2) IICA Morphisms:} every relation between elements is an
Idempotent, Immutable, Content-Addressed morphism; composition preserves these
properties, enabling arbitrarily nested world representations without new theory.
\textbf{(3) The Ashby--Bootstrap Correspondence:} Ashby's Requisite Variety Law,
interpreted naively, demands cardinality balance between connected systems. EWM resolves
this not by matching cardinalities directly, but through \textit{bootstrap
permutation}---multiplying measurements (e.g., n-grams) to bridge representational gaps
without sacrificing resolution. \textbf{(4) Noether Evolution:} structural change in an
EWM obeys conservation laws; every symmetry of the measurement regime yields a conserved
quantity. The lattice union invariant, the DRN decomposition (with its Lie-group
reading and N-steering control), the orthogonal channel decomposition, and the
holographic reconstruction property are all consequences of this principle. We
ground each principle
in the concrete instantiation of HLLSet Algebra and articulate the target formalism---a
Grothendieck topology on the category of measurement patches---toward which the framework
aspires, in spirit if not yet in letter.
\end{abstract}

% ==========================================================================
\section{Introduction}
% ==========================================================================

A world model is a system that learns an internal representation of its environment
from sensory input, then plans and reasons within that representation. Classical
approaches---JEPA, large language models, and token-based architectures---share a common
presupposition: the elements of the world (tokens, objects, features) are \textit{given}
before the model begins to operate. The vocabulary is fixed. The ontology is predefined.
The model's task is to learn the relations among these predefined elements.

\textbf{Emerging World Models} (EWM) invert this presupposition. In an EWM, elements are
not predefined. They \textit{emerge} from the accumulation of measurement interactions.
An element is not a token in a vocabulary; it is a stable pattern that crystallizes when
repeated measurements converge. Subsystems are not architectural modules; they are
coherent clusters of elements that co-occur across measurement trajectories.

This inversion has deep structural consequences. It shifts the foundation of world
modeling from \textit{representation} (encoding predefined elements into a vector space)
to \textit{measurement} (accumulating observations into a holographic plate from which
elements can be recovered). It replaces mutable, named parameters with immutable,
content-addressed bitmasks. It transforms learning from gradient descent over weight
matrices to rank rearrangement in a content-addressed dictionary.

This document consolidates the theoretical principles discovered across the HLLSet
Algebra research program \cite{astesj2026,sgs2024} into a unified structural framework.
We articulate four principles that govern every EWM, ground them in the concrete
instantiation of HLLSet Algebra, and identify Grothendieck topology as the target
formalism toward which the framework aspires.

\subsection{Why ``Emerging''?}

The term ``emerging'' is deliberate. It signals three properties that distinguish EWM
from classical world models:

\begin{enumerate}[leftmargin=*,nosep]
  \item \textbf{Ontological emergence.} The identity of an element is not assigned by a
    designer. It is \textit{recognized} by the system when a measurement pattern
    stabilizes across repeated observations. A token that appears once is noise. A token
    that persists across temporal layers becomes an element.
  \item \textbf{Structural emergence.} Relations between elements are not wired by an
    architect. They are \textit{computed} as IICA morphisms between the content-addressed
    fingerprints of the elements. The lattice of relations emerges from the lattice of
    measurements.
  \item \textbf{Evolutionary emergence.} The system does not converge to a fixed point.
    It continuously reorganizes as new measurements arrive, old patterns fade, and the
    balance between novelty and retention shifts. Evolution is structural---governed by
    conservation laws---rather than parametric.
\end{enumerate}

% ==========================================================================
\section{Principle 1: Emergent Ontology}
\label{sec:emergence}
% ==========================================================================

\begin{principle}[Emergent Ontology]
\label{prin:emergence}
The elements of an Emerging World Model are not defined a priori. They emerge from
the accumulation of measurement interactions. An element is \textit{what persists}
across repeated observation.
\end{principle}

\subsection{The Classical Approach: Predefined Ontology}

In classical world models, the ontology is fixed before the model encounters any data:
\begin{itemize}[leftmargin=*,nosep]
  \item \textbf{LLMs} operate on a fixed token vocabulary (50K--250K tokens) defined by
    a tokenizer (BPE, SentencePiece). The tokenizer is trained once and frozen. Every
    subsequent observation is projected into this fixed vocabulary.
  \item \textbf{JEPA} operates on patch embeddings derived from a pretrained encoder.
    The patches are predefined spatial or temporal segments.
  \item \textbf{Symbolic AI} operates on hand-crafted symbol tables and ontologies.
\end{itemize}

In all cases, the system cannot \textit{discover} a new element. A concept that falls
between tokens is represented as a combination of existing tokens, never as a primitive
in its own right. A novel pattern that does not align with the pretrained encoder's
latent dimensions is invisible.

\subsection{The EWM Approach: Measurement as Ontogenesis}

In an EWM, the only primitive is the \textbf{measurement event}: a token entering the
system, hashed to a bit position, accumulating a mark in the holographic plate. There is
no ``vocabulary''---only a Lookup Table (LUT) that tracks which tokens have been observed
at which positions, and how often.

An element emerges when a measurement pattern satisfies two conditions:
\begin{enumerate}[leftmargin=*,nosep]
  \item \textbf{Stability.} The same token or token cluster maps to consistent bit
    positions across multiple measurement events.
  \item \textbf{Persistence.} The token's Term Frequency (TF) at its bit positions
    remains above a threshold across temporal layers.
\end{enumerate}

A token that appears once and never again leaves a mark (its bit is set in the HLLSet),
but it does not become an element---its TF is negligible, its rank is low, and it does
not participate in relations. A token that appears thousands of times across weeks
accumulates high TF, high rank, and dense R-links to co-occurring tokens. It has
\textit{emerged} as an element.

\subsection{Formalization}

Let $\mathcal{T}$ be the space of all possible tokens (e.g., all byte strings). Let
$h: \mathcal{T} \to \mathcal{P}$ be an IICA morphism mapping tokens to patches (bit
positions) in a finite patch space $\mathcal{P} = \{0,\ldots,N-1\}$.

\begin{definition}[Measurement Event]
A measurement event at time $t$ is a pair $(t_i, p_i)$ where $t_i \in \mathcal{T}$ is a
token and $p_i = h(t_i) \in \mathcal{P}$ is its patch assignment.
\end{definition}

\begin{definition}[Element]
An element $e \in \mathcal{E}$ is a token $t \in \mathcal{T}$ whose accumulated
measurement count $\text{TF}(t)$ exceeds an emergence threshold $\theta_e$ and whose
temporal distribution across layers is non-uniform (it has a ``lifetime,'' not just a
single appearance):
\begin{equation}
  e \in \mathcal{E} \iff \sum_{i=0}^{K} \text{TF}_i(t) > \theta_e
  \;\land\; \exists i,j: \text{TF}_i(t) \gg \text{TF}_j(t)
\end{equation}
\end{definition}

The set of elements $\mathcal{E}$ is not fixed. It grows as new tokens accumulate
measurements, and it contracts as tokens fade. The ontology is \textit{emergent} and
\textit{dynamic}.

\subsection{HLLSet Instantiation}

In HLLSet Algebra, the emergent ontology is realized through the LUT (Lookup Table).
The LUT is a monotonic CRDT that maps each (register, trailing\_zeros) bucket to the
set of tokens that hash to that bucket, together with their accumulated TF values. The
LUT \textit{is} the system's ontology at any moment---the set of all tokens that have
been observed, with their measurement history.

A new token does not need to be ``added to the vocabulary.'' It is hashed, its bit is
set, and its entry in the LUT is created or incremented. The ontology grows
organically with every measurement event.

% ==========================================================================
\section{Principle 2: IICA Morphisms as Relational Foundation}
\label{sec:iica}
% ==========================================================================

\begin{principle}[IICA Relations]
\label{prin:iica}
Every relation between elements in an EWM is an \textbf{IICA morphism}: Idempotent,
Immutable, and Content-Addressed. Composition of IICA morphisms is IICA. This single
theorem eliminates the need for coordination, consensus, or naming at any scale of
system composition.
\end{principle}

\subsection{The Three Properties}

\begin{definition}[IICA Morphism]
A function $f: X \to Y$ is an IICA morphism iff:
\begin{itemize}[leftmargin=*,nosep]
  \item \textbf{Idempotent:} For every $x \in X$, computing $f(x)$ repeatedly---invoking
    $f$ on the same argument $x$ any number of times---always yields the same output
    $f(x)$.  When the codomain coincides with the domain ($X = Y$), this specializes to
    the classical algebraic form $f(f(x)) = f(x)$ (e.g., HLLSet union:
    $A \cup A = A$).  When $X \neq Y$ (e.g., a hash function
    $h: \text{Tokens} \to \{0,1\}^{160}$), the output type does not match the input type,
    so self-composition $h(h(x))$ is not meaningful; idempotence reduces to
    \textbf{determinism}: the same token $a$ hashes to the same digest $h(a)$ every time,
    regardless of when, where, or how many times $h$ is invoked.
  \item \textbf{Immutable:} $f(x) = y$ is fixed once computed. $y$ never changes.
  \item \textbf{Content-Addressed:} If $a = b$ then $f(a) = f(b)$. The output
    \textit{is} the address---given the content, you can find it again without a
    directory or a naming service.
\end{itemize}
\end{definition}

The canonical IICA morphism is a cryptographic hash function. $\text{SHA-1}(x)$ always
produces the same 160-bit digest for the same $x$ (idempotent as determinism); the digest
never changes (immutable); and the digest serves as the content address in IPFS/IPLD
(content-addressed).

\subsection{The Composition Theorem}

\begin{theorem}[IICA Composition]
\label{thm:composition}
Let $f_1, f_2, \ldots, f_n$ be IICA morphisms. Then their composition
$f_n \circ f_{n-1} \circ \cdots \circ f_1$ is an IICA morphism.
\end{theorem}

\begin{proof}
Each $f_i$ is idempotent, immutable, and content-addressed. Composition preserves each
property:
\begin{itemize}[leftmargin=*,nosep]
  \item \textbf{Idempotency:} For any input $x$, the composition
    $F = f_n \circ \cdots \circ f_1$ is a deterministic function: the same $x$ always
    triggers the same chain of intermediate results $f_1(x), f_2(f_1(x)), \ldots$ and
    produces the same final output $F(x)$.  When $X_1 = Y_n$ (the composition is an
    endomorphism), this additionally yields the algebraic form $F(F(x)) = F(x)$.
  \item \textbf{Immutability:} Each intermediate result $f_i(\ldots)$ is fixed once
    computed; the entire chain of intermediates is immutable.
  \item \textbf{Content-Addressability:} $a = b \implies f_1(a) = f_1(b) \implies
    \cdots \implies F(a) = F(b)$.
\end{itemize}
\end{proof}

This theorem is the engine of compositional emergence. It means that \textbf{any
pipeline of IICA morphisms is itself an IICA morphism}. You can compose tokenizers,
hashers, bridges between domains, materializers, and temporal aggregators without
ever introducing mutable state, naming conflicts, or coordination overhead.

\subsection{The IICA Pipeline}

In HLLSet Algebra, the full measurement pipeline is an IICA composition:

\begin{equation}
  \text{Tokens} \xrightarrow{h_1} \text{Token Hashes (u64)}
  \xrightarrow{h_2} \text{Bit Positions} \xrightarrow{h_3} \text{Bridge HLLSet}
  \xrightarrow{h_4} \text{Materialized Tokens} \xrightarrow{h_5} \cdots
\end{equation}

Each $h_i$ is a MurmurHash3-based IICA morphism. The composition $h_5 \circ h_4 \circ
h_3 \circ h_2 \circ h_1$ is IICA. This means:
\begin{itemize}[leftmargin=*,nosep]
  \item \textbf{No consensus protocols.} Two nodes ingesting the same tokens produce
    identical HLLSets. The lattice converges by construction.
  \item \textbf{No versioning.} HLLSets are immutable. ``Version $n+1$'' is a new
    HLLSet with a new CID. The old HLLSet remains valid and addressable.
  \item \textbf{No naming.} Every HLLSet's CID is its SHA-1 hash. No directory, no
    registry, no naming service.
  \item \textbf{Arbitrary nesting.} A swarm of $N$ robots, each with $M$ sensors,
    produces $N \times M$ IICA pipelines. Their composition at the command post is IICA.
    The theory does not change. The operations are identical. Only the scale differs.
\end{itemize}

\subsection{R-Links: IICA Morphisms as Objects}

A critical structural feature of EWM is that IICA morphisms between elements are
themselves elements. In HLLSet Algebra, the R-link $R_{AB} = A \cap B$ is both:
\begin{itemize}[leftmargin=*,nosep]
  \item A \textbf{morphism} from $A$ to $B$ (it captures the structural overlap that
    relates them).
  \item An \textbf{object} (it is an HLLSet with its own CID, storable in IPFS,
    composable with further intersections).
\end{itemize}

This self-referential property---morphisms as objects---is the structural signature of
a category that can model its own relations. It anticipates the Grothendieck topology
target: in a sheaf topos, the subobject classifier is itself an object of the topos.

% ==========================================================================
\section{Principle 3: The Ashby--Bootstrap Correspondence}
\label{sec:ashby}
% ==========================================================================

\begin{principle}[Ashby--Bootstrap Correspondence]
\label{prin:ashby}
Ashby's Requisite Variety Law implies that connected systems must match in
representational capacity. When cardinalities are mismatched, EWM resolves the
imbalance not by reducing capacity but by \textbf{bootstrap permutation}---multiplying
measurements to expand the effective cardinality of the smaller system until it
matches the larger.
\end{principle}

\subsection{Ashby's Law and the Cardinality Problem}

Ashby's Requisite Variety Law \cite{ashby1952} states that a controller must have at
least as much variety as the system it controls. In EWM terms: when two systems are
connected by IICA morphisms, their representational capacity must be balanced.

Consider the HLLSet architecture. The LUT (Lookup Table) maps tokens to bit positions.
For a vocabulary of 80K Chinese characters, the LUT has 80K entries. The HLLSet bitmask
has $1024 \times 32 = 32{,}768$ bit positions. A naive reading of Ashby's law would
suggest that the LUT and HLLSet are mismatched: 80K vs.\ 32K.

But this reading is superficial. The HLLSet does not store individual tokens---it
stores \textit{combinations} of tokens through hash collisions. The bitmask encodes the
OR of all tokens that hash to each position. The information capacity is not determined
by the number of bit positions but by the number of \textit{distinguishable bitmask
configurations}, which is $2^{32,768}$.

\subsection{Bootstrap Permutation}

The deeper problem is not total capacity but \textbf{measurement resolution}. When a
single token maps to a single bit position, the measurement is coarse: the bit is either
set or not, and the materializer cannot distinguish which of several tokens set it.

The solution is \textbf{bootstrap permutation}: multiply the measurements by applying a
family of permutations to the token stream before measurement.

\begin{definition}[Bootstrap Permutation]
Let $\Pi = \{\pi_1, \ldots, \pi_k\}$ be a family of permutations on token sequences.
Given a token stream $S = (t_1, \ldots, t_n)$, the bootstrapped measurement is:
\begin{equation}
  H_{\Pi}(S) = \bigoplus_{j=1}^{k} H(\pi_j(S))
\end{equation}
where $H(\pi_j(S))$ is the HLLSet of the permuted stream and $\oplus$ is bitwise OR.
\end{definition}

Each permutation $\pi_j$ produces a different ``view'' of the same token stream. The
union of these views creates a richer measurement---multiple bit positions per original
token, finer discrimination, and better materialization fidelity.

\begin{remark}[Omar Khayyam's Bootstrap]
This principle---precision through multiplication of primitive measurements---has a
striking historical precedent.  In 1079, Omar Khayyam reformed the Persian solar
calendar (the Jalali calendar) using only an astrolabe and basic naked-eye
observations---instruments orders of magnitude cruder than those available to
19th-century European astronomers.  Yet his calendar achieved an error of one day in
approximately 5,000 years, surpassing the Gregorian calendar's one day in 3,330
years.  The secret was not instrument quality but \textbf{measurement
multiplication}: Khayyam accumulated observations across many nights over many years,
and the statistical convergence of many primitive views produced precision that no
single observation could achieve.  Bootstrap permutation is the same idea in the
measurement-theoretic domain: the family of n-gram views replaces the family of
nightly observations, and the HLLSet union replaces the calendar average.
\end{remark}

\subsection{n-Grams as Bootstrap}

In the current HLLSet implementation, \textbf{n-grams} are the concrete bootstrap
family. Given a token stream $(t_1, t_2, t_3, t_4, \ldots)$, the 2-gram view is
$((t_1, t_2), (t_2, t_3), (t_3, t_4), \ldots)$ and the 3-gram view is
$((t_1, t_2, t_3), (t_2, t_3, t_4), \ldots)$. Each n-gram is hashed as a compound
token, producing a distinct set of bit positions.

\begin{proposition}[n-Gram Cardinality Expansion]
A vocabulary of $|V|$ tokens, measured with n-grams up to order $K$, produces up to
$\sum_{n=1}^{K} |V|^n$ distinct measurement points. For $|V| = 80{,}000$ and $K = 3$,
this is approximately $5.12 \times 10^{14}$ possible n-gram tokens, massively exceeding
the $32{,}768$ HLLSet bit positions.
\end{proposition}

The bootstrap permutation bridges the cardinality gap: a vocabulary that appears
``too small'' relative to the HLLSet is expanded through n-gram composition; a
vocabulary that appears ``too large'' is aggregated through hash collisions into the
fixed bit space. The balance is achieved through \textit{measurement multiplication},
not cardinality reduction.

\subsection{General Bootstrap Families}

n-grams are one family of bootstrap permutations. The principle generalizes:
\begin{itemize}[leftmargin=*,nosep]
  \item \textbf{Skip-grams:} $(t_i, t_{i+k})$ for various skip distances $k$.
  \item \textbf{Convolutional windows:} weighted combinations of neighboring tokens.
  \item \textbf{Spectral permutations:} frequency-domain rearrangements via FFT.
  \item \textbf{Random projections:} a family of random hash functions, each producing
    a distinct HLLSet view.
  \item \textbf{Cross-modal bootstraps:} the same stream measured through different
    sensory modalities, each with its own IICA pipeline.
\end{itemize}

The bootstrap family is a design parameter of the EWM. Different families produce
different emergent ontologies from the same underlying token stream. The choice of
family is the EWM analog of architecture selection in neural networks---it determines
what patterns the system can recognize.

% ==========================================================================
\section{Principle 4: Noether's Theorem as EWM Evolution Law}
\label{sec:noether}
% ==========================================================================

\begin{principle}[Noether Evolution]
\label{prin:noether}
In an EWM, every continuous symmetry of the measurement regime yields a conserved
quantity. Structural change is governed by conservation laws, not by gradient descent.
The system evolves by redistributing information across the lattice while preserving
global invariants.
\end{principle}

\subsection{From Physics to World Models}

In physics, Noether's theorem \cite{noether1918} establishes that every differentiable
symmetry of the action of a physical system corresponds to a conservation law:
time-translation symmetry $\rightarrow$ energy conservation; spatial-translation
symmetry $\rightarrow$ momentum conservation; rotational symmetry $\rightarrow$ angular
momentum conservation.

In an EWM, the ``action'' is the measurement regime---the accumulation of tokens into
HLLSets across temporal layers. The symmetries are transformations of the measurement
process that leave the structural relations invariant. The conserved quantities are
global properties of the lattice that persist across evolution steps.

\subsection{The Temporal Pyramid and [UM]-Net Cycles}

Before stating the invariant, we briefly introduce the two temporal structures that give
it physical meaning.

\textbf{The temporal pyramid} is a configurable $K$-layer compression stack (default
$K = 7$: second, minute, hour, day, week, month, year). Each layer $L_i$ accumulates
all measurement events within its temporal window via monotonic union:
$L_i = \bigcup S(t)$ for $t$ in the window. At each window boundary, the layer cascades
into the next coarser layer: $L_{i+1} \leftarrow L_{i+1} \cup L_i$, then $L_i$ resets.
After compression, layers are mutually exclusive---no time slice appears in more than one
layer. The complete system state is their union:
$H_{\text{system}} = \bigcup_{i=0}^{K} L_i$. Compression ratios range from $60\!:\!1$
(seconds$\to$minutes) to $12\!:\!1$ (months$\to$years); seven layers compress
approximately $3.15 \times 10^7$ seconds into one year-layer HLLSet.

\textbf{[UM] agents and [UM]-Net} provide a complementary temporal structure at shorter
timescales.  A \textbf{[UM]} (unified model) is the concrete implementation of a single
EWM: a processing node that receives a stream of materialized encodings from the external
environment \texttt{[E]}, ingests them into its internal HLLSet lattice, and materializes
the result back to a disambiguated encoding stream for downstream consumption.  Its
internal pipeline is a compound IICA morphism:

\begin{equation}
  \text{[UM]} \;=\;
  \text{(encodings)} \xrightarrow{\text{ingest}} \{\text{HLLSets}\}
  \xrightarrow{\text{materialize}} \text{(encodings)}
\end{equation}

The ingestion step hashes the incoming encodings into bit-vectors and accumulates them
into the lattice via monotonic union; the materialization step projects the lattice
through the agent's LUT and recovers ordered, disambiguated tokens via De Bruijn
reconstruction (Section~\ref{sec:hllset}).  Each individual step in the pipeline
(hashing, union, LUT projection) is IICA.  The [UM] node as a whole is IICA
\textit{between commits}: as long as its internal TF-Vec and HLLSet lattice are
fixed, the same input encodings always produce the same output encodings.  State
changes---new measurements accumulated, ranks updated---occur only at explicit
\textbf{commit events}, which define transactional boundaries.  Within a transaction,
[UM] behaves as a pure IICA morphism; across transactions, the same input may produce
different output as the lattice evolves.

A \textbf{[UM]-Net} is a directed graph $G = (A, E)$ of [UM] nodes connected by
fire-and-forget edges.  The external environment \texttt{[E]} feeds encodings into source
nodes; each [UM] processes independently and fires its output downstream with no
acknowledgment, no handshake, and no retry.  Convergence is guaranteed by IICA and CRDT
union at confluence nodes.  The full system diagram is:

\begin{equation}
  \texttt{[E]} \longrightarrow \text{(encodings)} \longrightarrow
  \texttt{[UM]}_1 \longrightarrow \texttt{[UM]}_2 \longrightarrow \cdots
\end{equation}

Critically, \textbf{cycles are not prohibited}: a recurrent loop
$a \to \cdots \to a$ in $G$ is a \textit{short-term memory} structure, resolved by
temporal separation---$\text{out}_a(t)$ feeds back as $\text{in}_a(t + k)$ after $k$
ingestion steps around the cycle.  Each pass through the cycle produces a new observation
at a later time, not a re-entry of the same data.  The evolution equation
$H(t) = H(S(t), H(t-1), D, R, N)$ captures both pyramid cascade and cycle recurrence
under a single dynamics.

\subsection{The Union Invariant}

The fundamental Noether symmetry in EWM is \textbf{path independence}: information can
travel through the temporal pyramid via multiple routes (L0$\to$L1$\to$L2, or
L0$\to$L2 directly, or L1$\to$L4$\to$L5). The conserved quantity is the union of all
temporal layers:

\begin{equation}
  H_{\text{system}}(t) = \bigcup_{i=0}^{K} L_i(t) = \text{constant over time}
\end{equation}

\begin{theorem}[Temporal Union Invariant]
\label{thm:union}
For any sequence of measurement events and any cascade policy, the bitwise union of all
temporal layers is monotonic and convergent. Information is never lost from
$H_{\text{system}}$; it is only redistributed across layers.
\end{theorem}

This is eventual consistency by construction---not as a protocol to implement, but as a
mathematical property of the lattice. No consensus algorithm. No leader election. No
retry logic. The lattice converges because the union is monotonic and multiple paths
guarantee that every bit reaches every layer that needs it.

\subsection{The DRN Conservation Law}

The DRN decomposition (Departed, Retained, Novel) expresses evolution as a conservation
balance:

\begin{equation}
  H(t) = H(S(t), H(t-1), D(t-1), R(t-1), N(t))
\end{equation}

where:
\begin{itemize}[leftmargin=*,nosep]
  \item $S(t)$ --- current observation as HLLSet
  \item $H(t-1)$ --- previous system state
  \item $D(t-1) = H(t-1) \setminus H(t)$ --- departed context
  \item $R(t-1) = H(t-1) \cap H(t)$ --- retained context
  \item $N(t) = H(t) \setminus H(t-1)$ --- novel context
\end{itemize}

The Noether steering condition is a conservation law on bit count across the DRN
transition:

\begin{equation}
  |\,\text{card}(N(t)) - \text{card}(D(t-1))\,| \to 0
\end{equation}

At steady state, the number of new bits entering the system equals the number of old
bits departing. The structural identity of the system is preserved even as its content
evolves. In rank-weighted form:

\begin{equation}
  \left|\sum_{b \in N(t)} R_b(t) - \sum_{b \in D(t-1)} R_b(t-1)\right| \to 0
\end{equation}

This is the EWM analog of energy conservation: the total ``rank mass'' of the system is
conserved across transitions. Learning is not accumulation---it is redistribution.

\subsection{The Lie-Group Reading of the DRN Decomposition}
\label{sec:lie-drn}

The DRN decomposition has a precise group-theoretic interpretation connecting
it to on-manifold optimization for rigid-body motion \cite{argueta2026}. The
patch space $\mathcal{B} = \{0,1\}^{32{,}768}$ is not only a join-semilattice
under OR; it is a \textbf{Boolean ring}, hence an abelian group under symmetric
difference (XOR). For a transition $A = H(t-1) \to B = H(t)$:

\begin{equation}
  S \;=\; A \oplus B \;=\; (A \setminus B) \sqcup (B \setminus A)
  \;=\; D(t-1) \sqcup N(t),
\end{equation}

where $\sqcup$ denotes disjoint union. The retained context is the fixed part
$R(t-1) = A \cap B$, and the two states decompose as

\begin{equation}
  A = R \oplus D, \qquad B = R \oplus N, \qquad B = A \oplus S.
\end{equation}

Thus $(N, D)$ are the coordinates of the transition in the tangent space at
$A$---the DRN split is the ``logarithm'' of the group relative transform, with
$R$ as the stabilizer. The cardinality balance follows immediately:

\begin{equation}
  |B| = |A| - |D| + |N|,
\end{equation}

so the Noether steering condition
$|\,\text{card}(N(t)) - \text{card}(D(t-1))\,| \to 0$ is equivalent to
$|H(t)| = |H(t-1)|$. Geometrically, the update has \textbf{zero radial
component}: the state's velocity is tangent to the Hamming sphere of constant
cardinality---a pure ``rotation'' of structure, in which content changes while
the measure is preserved. In the language of on-manifold optimization,
stability as $D$--$N$ balance is not a flat-space approximation but the correct
first-order on-manifold condition.

\begin{remark}[Gauge Fixing: Where $D$ and $N$ Must Be Measured]
The Temporal Union Invariant (Theorem~\ref{thm:union}) makes
$H_{\text{system}}$ monotone: bits never leave the system, so at the system
level $D \equiv 0$ and the balance can never hold. Conversely, at the raw layer
level, the cascade moves a bit from $L_0$ to $L_1$, recording a departure in one
layer and a novelty in the other---a ``wrap-around'' artifact analogous to
subtracting the angles $179^\circ$ and $-179^\circ$ and obtaining $358^\circ$
instead of $2^\circ$. The DRN balance must therefore be measured on the
\textbf{operative state}---the thresholded working set $\mathcal{S}(t)$ of
Section~\ref{sec:cortex-learning}---modulo the cascade symmetry. Departure
means \textit{stopped attending}, not erased: an excluded HLLSet remains
content-addressed and re-admissible at any time.
\end{remark}

\subsection{The Holographic Property as Noether Consequence}

The holographic property---that the lattice top $H_{\text{system}}$ plus the TF stack
can reconstruct any past state---is a direct consequence of the union invariant:

\begin{equation}
  \text{past\_state}(t) \approx H_{\text{system}}(\text{now}) \odot
  \text{TF}_{\text{stack}}[t]
\end{equation}

Because $H_{\text{system}}$ never loses information (union invariant), and the TF stack
records \textit{when} each bit position was active (temporal distribution), the
reconstruction is always possible. The information is not stored separately for each
time step---it is encoded holographically in the lattice, with the TF stack serving as
the reference beam for temporal selectivity.

\subsection{The Fisher Matrix as Structural Derivative}

The cross-layer Fisher-like matrix captures the second-order structure of evolution:

\begin{equation}
  F_{bb'} = \sum_{i=0}^{K} B^{(i)}_b \cdot B^{(i)}_{b'}
\end{equation}

where $B^{(i)}_b = 1$ if bit $b$ is set in layer $i$. $F_{bb'}$ counts how many layers
have both bits $b$ and $b'$ set simultaneously. High $F_{bb'}$ indicates systemic
co-occurrence (a structural invariant). Low $F_{bb'}$ indicates noise. The Fisher matrix
enables the Noether controller to distinguish isolated fluctuations from systemic phase
transitions.

\subsection{Noether Controller}

The Noether controller monitors the cross-layer relationship matrix and decides which
temporal layer should drive the next action. When the BSS between L0 (current second)
and L1 (recent minute) drops below a threshold, the controller detects divergence and
overrides instinct with broader context. When L0 diverges from L3 (day-scale goals),
the controller triggers re-routing.

The controller is the \textbf{structural actuator} of the Noether principle. It does not
adjust weights---it selects which layer of accumulated measurement history is relevant
to the current situation. The six control parameters (pyramid depth, rank functions,
attention threshold, steering thresholds) are the knobs of structural evolution, not the
parameters of a learned function.

\subsection{N-Steering: Novelty as Control Input}
\label{sec:n-steering}

In on-manifold optimization one chooses a full tangent vector and retracts it
onto the manifold. The EWM controller, by contrast, directly controls only the
\textbf{positive} component of the transition---the novelty channel $N$---via
the gate thresholds $\theta_{\tau}, \theta_{\rho}$, the bootstrap depth, and the
measurement policy. The departure $D$ is the passive response (window expiry,
rank decay, cascade). N-steering is therefore a \textbf{one-sided control
problem}: steer the innovation axis so that it tracks the predicted decay.

\begin{definition}[N-Steering Law]
Let $\widehat{D}(t)$ be the departure predicted by a Type-1 forward model from
the current lattice state and control settings. The N-steering law sets the
novelty throttle so that
\begin{equation}
  \mathbb{E}\bigl[\text{card}(N(t))\bigr]
  \;\to\; \mathbb{E}\bigl[\text{card}(\widehat{D}(t))\bigr],
\end{equation}
enforcing zero \textit{expected} radial drift:
$\mathbb{E}[|H(t)|] = \mathbb{E}[|H(t-1)|]$. The Noether condition becomes a
control objective, not a post-hoc diagnostic.
\end{definition}

The two-stage form of N-steering follows from the measurement structure of
Section~\ref{sec:channels}. First \textbf{unwrap}: a novelty spike in the coarse
channel may be the same structure measured at insufficient resolution; deepening
the bootstrap re-expresses apparent novelty as retained context before any
balance is attempted. Second \textbf{balance}: only the residual novelty---that
which survives re-measurement at higher n-gram order---is balanced against the
predicted departure. This is the EWM analog of the on-manifold rule: do not
subtract raw angles, unwrap first.

\subsection{Noether Steering Validation: The Measurement Hierarchy}
\label{sec:newton-hierarchy}

The N-steering law of Section~\ref{sec:n-steering} is validated by three
measurements of the transition, read directly from the immutable lattice. None
of the three modifies $S(t)$, $H(t)$, $H(t-1)$, or the DRN components---the
HLLSets are unchanged after measurement. None creates new bit positions
($\langle \text{reg}, \text{zeros} \rangle$ pairs): new bits in an HLLSet
vector can be created only during ingestion. The three orders are therefore
\textbf{non-invasive observables}: each is a pure read of the lattice, and no
order is obtained by differentiating a lower-order estimate. The hierarchy is
one of measurement, not of error propagation.

Define, on the causal pair $(N(t), D(t-1))$,

\begin{align}
  v(t) &= \text{card}(N(t)) - \text{card}(D(t-1)), \\
  \delta N(t) &= \text{card}(N(t)) - \text{card}(N(t-1)), \\
  \delta D(t) &= \text{card}(D(t-1)) - \text{card}(D(t-2)), \\
  a(t) &= \delta N(t) - \delta D(t) \;=\; v(t) - v(t-1), \\
  F(t) &= m(t)\, a(t),
\end{align}

where $m(t)$ is the rank mass of $S(t)$ (the raw cardinality
$\text{card}(S(t))$ is the cruder special case) and
$D(t-2) = H(t-2) \setminus H(t-1)$ is the departure of the preceding
transition. The three validation measurements are:

\begin{enumerate}[leftmargin=*,nosep]
  \item \textbf{Zeroth order (homeostasis):}
    $|\,\text{card}(N(t)) - \text{card}(D(t-1))\,| \to 0$---the original
    Noether steering condition, equivalently $v(t) = 0$.
  \item \textbf{First order (inertial stability):} $\delta N(t) = \delta D(t)$,
    equivalently $a(t) = 0$: the net flux is constant, so the system may stand
    still or grow or decay at a steady rate. Uniform flux---not flux equal to
    zero---is the inertial signature of stability.
  \item \textbf{Second order (steering force):} $F(t) = m(t)\, a(t)$---the
    force the N-steering controller must apply to change the trajectory.
\end{enumerate}

Each order is an independent measurement of the same underlying state: the
first-order condition does not inherit estimation error from the zeroth, and
the second does not inherit from the first, because all are read directly from
the lattice. The three orders form a jet hierarchy $J^0 \subset J^1 \subset
J^2$ over the state trajectory.

\begin{remark}[Newtonian Reading as Analogy]
The three orders read as the three laws of Newtonian dynamics. The analogy is
intended to identify the roots of the thinking, not to install dynamics on the
lattice:

\begin{itemize}[leftmargin=*,nosep]
  \item \textbf{First law.} $a(t) = 0$ implies $v(t) = \text{const}$: the
    system stays still ($v = 0$, the zeroth-order condition) or moves with
    steady flux (uniform growth or decay). In scalar form this is uniform
    motion of the measure; ``straight-line'' motion in the full Boolean ring
    additionally requires the direction of the transition
    $S(t) = N(t) \sqcup D(t-1)$ to be stationary.
  \item \textbf{Second law.} The steering force is the product $F = m\,a$,
    not the pair $[m, a]$. The channel accelerations combine signed:
    $a = \delta N - \delta D$. Taking magnitudes per channel before
    subtraction would erase the force precisely in the regime where ingress
    and egress diverge in opposite directions.
  \item \textbf{Third law.} Action and reaction at the $D$--$N$ boundary follow
    not from the second law but from momentum conservation---the union
    invariant (Theorem~\ref{thm:union}) and the rank-mass conservation of
    Principle~4. The internal channel forces cancel; only the net force
    $m\,a$ drives the center of mass.
\end{itemize}
\end{remark}

Because the measurements are exact reads of the immutable lattice, the
remaining variability in $a(t)$ and $F(t)$ is environmental, not measurement
error. The Noether controller may smooth the force signal over the temporal
pyramid before acting; this is control-theoretic filtering, not a correction
of the measurement.

% ==========================================================================
\section{HLLSet Algebra: Concrete Instantiation}
\label{sec:hllset}
% ==========================================================================

The four principles above are not abstract. They are realized in the HLLSet Algebra
framework \cite{astesj2026}, whose components map directly onto the EWM structure.

\subsection{The Patch Space as Measurement Site}

The HLLSet bitmask is a \textbf{holographic plate} of $32{,}768$ patches (bit
positions), each recording a binary mark: 0 = untouched, 1 = at least one token hashed
here. The hash function $h: \mathcal{T} \to \{0,\ldots,32767\}$ is the IICA measurement
morphism. The LUT $\mathcal{L}$ tracks the inverse image of $h$, enabling materialization
(recovering which tokens could have set which bits).

This is the minimal EWM: a finite set of measurement sites, an IICA morphism from tokens
to patches, and a monotonic accumulation operation (bitwise OR) that makes each patch a
join-semilattice. Principle~1 (emergent ontology) is realized by the LUT. Principle~2
(IICA relations) is realized by the hash pipeline and R-link lattice operations.

\subsection{The Temporal Pyramid as Noether Structure}

The 7-layer temporal pyramid (second, minute, hour, day, week, month, year) implements
Principle~4 (Noether evolution). Each layer accumulates measurements over its temporal
window via union. The cascade policy (L0 overflows into L1, L1 into L2, etc.) is a
discrete-time dynamical system whose conserved quantity is the union invariant.

\subsection{n-Grams as Bootstrap}

The tokenizer's n-gram pipeline (1-grams, 2-grams, 3-grams with boundary padding)
implements Principle~3 (bootstrap permutation). Each n-gram order is a distinct
measurement view of the same token stream. The union of n-gram HLLSets produces a richer
representation than any single order alone.

\subsection{Orthogonal Channel Decomposition and Error Narrowing}
\label{sec:channels}

The n-gram bootstrap has a structural refinement: the 1-gram, 2-gram, and
3-gram views are hashed into \textbf{pairwise disjoint} bit ranges. Within
every HLLSet, $H_1 \cap H_2 = H_1 \cap H_3 = H_2 \cap H_3 = \emptyset$; each
HLLSet is therefore a disjoint union

\begin{equation}
  H = H_1 \sqcup H_2 \sqcup H_3, \qquad |H| = |H_1| + |H_2| + |H_3|,
\end{equation}

where $H_n = H \cap G_n$ and $G_n$ is the global order-$n$ gate (bit-range
mask). The three channels are \textbf{orthogonal} components of the
measurement: a dense 3-gram population cannot pollute the 1-gram channel, and
vice versa.

\begin{proposition}[Additivity of BSS over Channels]
\label{prop:channel-bss}
For HLLSets $A = A_1 \sqcup A_2 \sqcup A_3$ and
$B = B_1 \sqcup B_2 \sqcup B_3$ with pairwise disjoint channels, all
cross-terms vanish and
\begin{equation}
  A \cap B = (A_1 \cap B_1) \sqcup (A_2 \cap B_2) \sqcup (A_3 \cap B_3),
\end{equation}
so the aggregate BSS is the convex combination
\begin{equation}
  \text{BSS}_{\tau}(A \to B) = \sum_{n=1}^{3} w_n\,
  \text{BSS}_{\tau}(A_n \to B_n), \qquad w_n = \frac{|B_n|}{|B|},
\end{equation}
and likewise for $\text{BSS}_{\rho}$. The DRN components are additive:
$N = N_1 \sqcup N_2 \sqcup N_3$, $D = D_1 \sqcup D_2 \sqcup D_3$, and
\begin{equation}
  |N| - |D| = \sum_{n=1}^{3} \bigl(|N_n| - |D_n|\bigr).
\end{equation}
\end{proposition}

\begin{proof}
For $n \neq m$, $A_n \cap B_m \subseteq G_n \cap G_m = \emptyset$. All
intersection and difference cardinalities therefore decompose additively over
the channels.
\end{proof}

\begin{remark}[Token-Level Nesting vs.\ Bit-Level Orthogonality]
The channels are orthogonal in the \textit{image} (bit space), while the token
structures they measure are nested in the \textit{preimage}: a shared 3-gram
contains a shared 2-gram and shared 1-grams, so at the token level relational
overlap implies element overlap---but not conversely (a shared element in
different contexts has empty 2- and 3-gram intersections). The nesting
information lives in the LUT; the bitmasks provide three independent
coordinates, and the LUT supplies the semantic relations between them.
\end{remark}

\subsubsection{Error Narrowing}

The orthogonality of the channels makes the three channel-wise BSS values
independent observables, narrowing measurement error in two ways.

First, each channel yields an \textbf{exact disjointness certificate}. Because
hashing is deterministic (IICA), a shared n-gram always sets the same bits in
both measurements; contrapositively,

\begin{equation}
  H_n(A) \cap H_n(B) = \emptyset \;\Longrightarrow\;
  \text{$A$ and $B$ share no $n$-grams} \qquad \text{(exactly)},
\end{equation}

with zero false negatives: the only error mode of the bitwise intersection is
false positives (hash collisions between distinct n-grams).

Second, the false-positive error composes cleanly. For disjoint channel
occupancies $o_{A,n}, o_{B,n}$ in a channel of $m_n$ bits, the expected fake
overlap is $o_{A,n}\, o_{B,n} / m_n$. The aggregate fake overlap is the sum
$\sum_n o_{A,n}\, o_{B,n} / m_n$, while the \textbf{joint agreement test}---
requiring non-empty intersection in all three channels before declaring genuine
overlap---has error $\varepsilon_1 \varepsilon_2 \varepsilon_3$: because the
channels are disjoint, collision events are independent across orders and the
error narrows multiplicatively.

\begin{remark}[Bit-Budget Trade-off]
Orthogonality costs resolution: each channel receives only a fraction of the
$32{,}768$ bits, so per-channel collision rates are higher than in a shared
space. The near-optimal allocation equalizes density, $m_n$ proportional to the
expected occupancy of order $n$. For streaming measurement the three
occupancies are comparable and an equal split is adequate; for corpus-scale
3-gram accumulation the 3-gram channel saturates first and deserves a larger
range.
\end{remark}

\subsubsection{Materialization as Local Repair}

The 3-gram channel is nearly deterministic: for almost every set bit
$b \in H_1$, the 3-gram evidence reduces the candidate set $\mathcal{L}_1[b]$
(the LUT entries hashing to $b$) to a singleton. Define the partial section

\begin{equation}
  \phi(b) = \text{the unique } x \in \mathcal{L}_1[b]
  \text{ whose 3-gram neighborhood is present in } H_3.
\end{equation}

Where $\phi$ is single-valued, materialization is the evaluation of a
section---no global reconstruction. Only on the ambiguity set
$\{b : |\phi^{-1}(b)| > 1\}$, of collision measure $\sim 10^{-5}$, is a local
repair needed (TF ranking or local De Bruijn). The materializer is thereby
demoted from a pipeline stage to an exception handler; the lattice dynamics
(BSS, DRN, N-steering) never require it. Materialization remains necessary only
at the I/O boundary, where the system must emit encodings to an external
decoder.

\subsection{Rank Algebra as Learning}

Learning is rank rearrangement. HLLSets are immutable (IICA); only their
\textit{relevance} to the current context changes. This section explains how ranks
are computed, why Forth hosts the ranking, and how the five-level propagation
distinguishes bit-level signal from HLLSet-level importance.

\subsubsection{TF vs.\ Rank: The Fundamental Separation}

\textbf{TF is stored. Rank is derived. They are not the same thing.}  The Term
Frequency vector (TF-Vec) is attached to to each EWM LUT(s), it's shared, monotonic CRDT of $32{,}768$ f64 values---one
per bit position---updated only during ingestion.  Rank is computed locally from TF
at query time and is never persisted in the protocol.  This separation has two
consequences.  First, TF requires the same token base for comparison (it is
hash-function-dependent); rank is hash-function-agnostic and meaningful across
different domains connected by the Universal Bridge.  Second, a HLLSet created years
ago can rise in relevance today because the TF-Vec evolved around it---new
measurements make its bit positions hot; stale positions quietly cool.  The HLLSet
itself never changed.

\subsubsection{The Forth Dictionary as Learning Engine}

The Forth DSL serves as the \textbf{canonical AST} that unifies all backends: the
same Forth source compiles identically to Lua (software simulation), Rust (native
execution), and Verilog (FPGA synthesis).  The Forth dictionary associates every
defined word---including every HLLSet---with a \textbf{rank}.  The dictionary is the
only component that learns:

\begin{center}
\begin{tabular}{@{}lll@{}}
\toprule
\textbf{Component} & \textbf{Operation} & \textbf{Learns?} \\
\midrule
Tokenizer       & bytes $\to$ HLLSet           & No (deterministic function) \\
Materializer    & collects HLLSet outputs       & No (passive observer) \\
\textbf{Forth Dictionary} & \textbf{assigns ranks} & \textbf{Yes (only this)} \\
\bottomrule
\end{tabular}
\end{center}

A rank adjustment is a structural change---which HLLSet drives action---without any
parametric change to the HLLSets themselves.  This satisfies Principle~4: learning is
redistribution of attention across the lattice, not accumulation of weight updates.

\subsubsection{Five-Level Rank Propagation}

Rank propagates from raw token frequency to compound lattice operations through five
compositional levels.  Each level aggregates the rank signal from the level below via
a pluggable function (F, G, H, K, L, M); the architecture specifies the framework,
the application chooses the functions.  All aggregation functions must be monotonic
to preserve CRDT convergence.

\begin{center}
\begin{tabular}{@{}ccl@{}}
\toprule
\textbf{Level} & \textbf{Function} & \textbf{Definition} \\
\midrule
5. Compound HLLSet rank & $L, M$ & $L(\max\{R_i\})$ for union,
                                    $M(\min\{R_i\})$ for intersection \\
4. HLLSet rank           & $K$     & $K(\text{degree in lattice DAG})$ \\
3. Register rank         & $H$     & $H(\{\text{bit-R}[tz] \mid tz \in 0..31\})$ \\
2. Bit rank              & $G$     & $G(\{\text{token-R} \mid \text{all tokens hashing to }(r, tz)\})$ \\
1. Token rank            & $F$     & $F(\text{TF})$ --- rank from token frequency \\
\bottomrule
\end{tabular}
\end{center}

\textbf{Levels 1--3} operate within a single HLLSet: token frequency $\to$ bit
position activity $\to$ register importance.  \textbf{Level 4} lifts to the lattice
DAG: a HLLSet's structural importance depends on how many R-links connect it to other
HLLSets (its degree in the intersection graph).  \textbf{Level 5} composes ranks
across lattice operations: the rank of a union is the maximum of its constituents'
ranks; the rank of an intersection is the minimum.

This hierarchy means that a bit position with high token TF (Level~1) propagates
signal upward, eventually making the HLLSets that contain it (Level~4) and the
compound operations that include them (Level~5) more likely to drive behavior.  The
Forth dictionary's rank assignment is the final projection of this five-level signal
onto the action-selection mechanism.

% ==========================================================================
\section{Toward Grothendieck Topology}
\label{sec:gt}
% ==========================================================================

The four principles above point toward a deeper mathematical structure:
\textbf{Grothendieck topology} \cite{grothendieck1972} on the category of measurement
patches. We articulate this target not as a completed formalism, but as a compass
bearing---the ``spirit'' in which the EWM framework is developed.

\subsection{The Site: A Category of Patches}

Let $\mathcal{P}$ be a finite discrete category whose objects are patch addresses and
whose only morphisms are identities. This is the \textbf{site}---the space of measurement
locations.

A \textbf{presheaf} on $\mathcal{P}$ assigns to each patch $p$ a set of possible marks
$\Sigma(p)$ (the mark alphabet) and to each morphism (only identities) the trivial
restriction. An HLLSet is a global section of this presheaf:

\begin{equation}
  H \in \Gamma(\mathcal{P}, \Sigma) = \prod_{p \in \mathcal{P}} \Sigma(p)
\end{equation}

For standard HLLSets, $\Sigma(p) = \{0, 1\}$ (a bit) with accumulation $\oplus =
\text{OR}$. For the generalized case (GENERAL\_HLLSET\_ALGEBRA.md), $\Sigma(p)$ can be
any finite join-semilattice: $\{0, 1, 2\}$ (trits with max), $\{0, 1, 2, 3\}$ (quits),
or arbitrary discrete alphabets.

\subsection{Coverages and the Gluing Condition}

A Grothendieck topology on $\mathcal{P}$ specifies which families of morphisms
$\{U_i \to U\}$ count as \textbf{covering families}. The gluing condition states: if
data on each $U_i$ is consistent on overlaps $U_i \times_U U_j$, then there exists
unique data on $U$ that restricts to the given data on each $U_i$.

In EWM terms, the covering families are the \textbf{bootstrap permutation family}
(Principle~3). Each permutation $\pi_j$ provides a partial view of the token stream
(a measurement on a ``sub-site''). The gluing condition is satisfied when these partial
views are consistent---when the same token maps to compatible marks across different
permutations.

The HLLSet union $\bigcup_{\pi \in \Pi} H_\pi$ is the \textbf{sheafification} of the
bootstrap family: it combines all partial measurements into a global section.

\subsection{Geometric Morphisms as Universal Bridges}

A geometric morphism $f: \text{Sh}(\mathcal{P}) \to \text{Sh}(\mathcal{Q})$ between
sheaf toposes on different patch spaces is the categorical analog of the Universal
Bridge. It consists of:
\begin{itemize}[leftmargin=*,nosep]
  \item A \textbf{direct image} $f_*$ that pushes HLLSets from $\mathcal{P}$ to
    $\mathcal{Q}$.
  \item An \textbf{inverse image} $f^*$ that pulls HLLSets from $\mathcal{Q}$ back to
    $\mathcal{P}$.
\end{itemize}

In the current implementation, the Universal Bridge implements the direct image via
re-representation: format each bit position as a token, re-hash into the target space.
The inverse image (materialization from target back to source) is approximate---it
depends on the LUT and is subject to the Statistics Constraint (TF vectors are not
transferred across bridges).

\subsection{The Subobject Classifier and R-Links}

In a topos, the \textbf{subobject classifier} $\Omega$ classifies subobjects: for any
object $X$, there is a bijection between subobjects of $X$ and morphisms $X \to \Omega$.
For the Boolean topos on $\mathcal{P}$, $\Omega$ is the two-element set $\{0, 1\}$
with the OR accumulation. Every HLLSet $H$ is a morphism $H: \mathcal{P} \to \Omega$.

The R-link $R_{AB} = A \cap B$ is the \textbf{pullback} of $A$ and $B$ over $\Omega$.
In categorical terms:
\begin{equation}
  \begin{CD}
    R_{AB} @>>> A \\
    @VVV @VVV \\
    B @>>> \Omega
  \end{CD}
\end{equation}

This is why R-links are both morphisms and objects: they are pullbacks in the topos.
The lattice of HLLSets is a \textbf{Heyting algebra}---the internal logic of the topos.

\subsection{What the Grothendieck Target Gives Us}

Adopting Grothendieck topology as the target formalism would provide:

\begin{enumerate}[leftmargin=*,nosep]
  \item \textbf{A rigorous gluing theorem.} Given consistent measurements on a covering
    family, the global state is uniquely determined. This formalizes Principle~3 (why
    bootstrap permutations are sufficient).
  \item \textbf{A systematic theory of bridges.} Geometric morphisms between patch
    topoi classify all possible Universal Bridges. Different choices of direct/inverse
    image correspond to different bridge implementations.
  \item \textbf{An internal logic.} The Heyting algebra structure on the HLLSet lattice
    is the internal logic of the topos. Union, intersection, and implication correspond
    to logical OR, AND, and IMPLIES. The EWM \textit{reasons} within its own lattice.
  \item \textbf{Cohomology.} Sheaf cohomology groups would measure the ``obstruction''
    to gluing --- quantifying how much information is lost when different bootstrap views
    disagree.
\end{enumerate}

We are not yet ready to formalize EWM as a full Grothendieck topology. The current
framework is a \textbf{pre-sheaf structure} on a discrete site with a bootstrap covering
family. The four principles---emergent ontology, IICA morphisms, bootstrap permutation,
and Noether evolution---are the structural requirements that any such topology must
satisfy. The destination may differ from the current compass bearing. But the bearing
itself---gluing local measurements into global states through consistent covering
families---is the ``spirit'' that guides the framework.

% ==========================================================================
\section{Discussion}
\label{sec:discussion}
% ==========================================================================

\subsection{Unification}

The four principles unify the existing body of work:

\begin{itemize}[leftmargin=*,nosep]
  \item \textbf{HLLSet Theory} \cite{astesj2026} establishes the base lattice
    (Principle~1, 2) and the IICA gate (Principle~2).
  \item \textbf{The Self-Reprogramming Architecture} and \textbf{STANDARD.md} specify
    the temporal pyramid, DRN decomposition, Noether controller, and holographic memory
    (Principle~4).
  \item \textbf{GENERAL\_HLLSET\_ALGEBRA.md} identifies the categorical structure
    (Bool$_{32768}$, IICA\_Hash, Patch category) and the Grothendieck target
    (Section~\ref{sec:gt}).
  \item \textbf{The n-gram tokenizer} and \textbf{Universal Bridge} implement
    bootstrap permutation (Principle~3).
  \item \textbf{CAAL-LLM} validates that the principles produce working intelligence:
    80\% accuracy on driving scenarios from 10 sentences, no gradient descent
    \cite{astesj2026}.
\end{itemize}

\subsection{What EWM Is Not}

It is important to state what EWM is \textit{not}:

\begin{itemize}[leftmargin=*,nosep]
  \item \textbf{Not a neural network.} There are no weight matrices, no activation
    functions, no backpropagation. The lattice is the computation; ranks are the
    attention; the dictionary is the memory.
  \item \textbf{Not a token-based model.} There is no fixed vocabulary. Elements emerge
    from measurement accumulation. A new token is hashed and enters the lattice without
    architectural change.
  \item \textbf{Not a probabilistic model in the Bayesian sense.} There are no
    priors, no posteriors, no sampling. The BSS (Bell State Similarity) is a
    deterministic bitwise measure, not a probability.
  \item \textbf{Not a database.} Although content-addressed and IPFS-native, the
    lattice is a computational structure, not a storage system. The HLPP protocol
    specifies persistence, but the lattice operations are the computation.
\end{itemize}

\subsection{Open Questions}

\begin{enumerate}[leftmargin=*,nosep]
  \item \textbf{Optimal bootstrap families.} What is the minimal set of permutations
    that guarantees $\epsilon$-accurate materialization for a given token stream
    distribution? The current n-gram choice is empirical.
  \item \textbf{Convergence rates.} Under what conditions does the Noether controller
    converge to a stable layer selection policy? What is the relationship between
    environmental volatility and convergence time?
  \item \textbf{Topos formalization.} What is the minimal Grothendieck topology that
    captures the bootstrap covering condition? Is the current discrete site sufficient,
    or do we need a non-trivial coverage?
  \item \textbf{Cross-modal emergence.} If tokens come from different sensory modalities
    (text, image patches, audio spectra), how does the unified lattice represent
    cross-modal elements? The Universal Bridge provides translation; emergence across
    modalities is unexplored.
  \item \textbf{The Statistics Constraint.} Why must TF vectors NOT be transferred
    across bridges? The constraint is empirical (LUT initialization failure). A
    categorical explanation---TF as a sheaf of local statistics that does not push
    forward under geometric morphisms---is plausible but unproven.
\end{enumerate}

\subsection{The Road Ahead}

The EWM framework is at the beginning of its development. The four principles provide
a structural compass. The HLLSet Algebra implementation provides a concrete testbed.
The Grothendieck target provides theoretical direction. The path forward involves:

\begin{enumerate}[leftmargin=*,nosep]
  \item \textbf{Harden the implementation.} Complete the FPGA backend, distributed mesh
    networking, and self-ingestion pipeline to validate the principles at scale.
  \item \textbf{Explore the bootstrap design space.} Systematic experiments with
    different permutation families, mark alphabets (trits, quits), and accumulation
    semilattices.
  \item \textbf{Develop the categorical theory.} Move from pre-sheaf to sheaf, from
    discrete site to Grothendieck topology, from Bridge to geometric morphism.
  \item \textbf{Bridge to applications.} The CAAL-LLM proof-of-concept demonstrates
    viability. Real-time adaptive systems (robotics, autonomous vehicles, industrial
    control) are the natural deployment targets for an FPGA-native, gradient-free world
    model.
\end{enumerate}

% ==========================================================================
% ── References ────────────────────────────────────────────────────────────
% ==========================================================================

\begin{thebibliography}{9}

\bibitem{astesj2026}
A.~Mylnikov.
\emph{HLLSet Theory: A Unified Framework for Probabilistic Knowledge Representation}.
ASTESJ, Vol.~11, No.~2, pp.~62--74, 2026.
\url{https://www.astesj.com/publications/ASTESJ_110202.pdf}

\bibitem{sgs2024}
A.~Mylnikov.
\emph{Self Generative Systems (SGS) and Its Integration with AI Models}.
AISNS '24: Proceedings of the 2024 2nd International Conference on Artificial
Intelligence, Systems and Network Security, pp.~345--354, 2024.
\url{https://doi.org/10.1145/3714334.3714392}

\bibitem{ashby1952}
W.~R.~Ashby.
\emph{Design for a Brain}.
Chapman \& Hall, 1952.

\bibitem{noether1918}
E.~Noether.
\emph{Invariante Variationsprobleme}.
Nachr.~d.~K\"onig.~Gesellsch.~d.~Wiss.~zu G\"ottingen, Math-phys.~Klasse,
pp.~235--257, 1918.

\bibitem{flajolet2007}
P.~Flajolet, \'{E}.~Fusy, O.~Gandouet, F.~Meunier.
\emph{HyperLogLog: the analysis of a near-optimal cardinality estimation algorithm}.
AofA, 2007.

\bibitem{lecun2022}
Y.~LeCun.
\emph{A Path Towards Autonomous Machine Intelligence}.
Open Review, 2022.
\url{https://openreview.net/forum?id=BZ5a1r-kVsf}

\bibitem{grothendieck1972}
A.~Grothendieck, J.~L.~Verdier.
\emph{Th\'eorie des Topos et Cohomologie \'Etale des Sch\'emas} (SGA 4).
Lecture Notes in Mathematics, Vols.~269, 270, 305. Springer, 1972--1973.

\bibitem{maclane1992}
S.~Mac Lane, I.~Moerdijk.
\emph{Sheaves in Geometry and Logic: A First Introduction to Topos Theory}.
Springer, 1992.

\bibitem{argueta2026}
C.~Argueta.
\emph{Lie Groups and On-Manifold Optimization}.
Soul Hackers Labs, Medium, July 2026.
\url{https://soulhackerslabs.com/lie-groups-and-on-manifold-optimization-b0814ae7511b}

\bibitem{lincoln2026}
A.~Lincoln.
\emph{Uncertainty reduction as a candidate primary reinforcer: an evolutionary
and neural account}.
Academia Neuroscience and Brain Research, Vol.~2, 2026.
\url{https://doi.org/10.20935/AcadNeurosci8451}

\end{thebibliography}

% == Acknowledgment =========================================================
\section*{Acknowledgment}
This consolidation was developed in dialogue with DeepSeek Code, an AI research
assistant. The four principles---Emergent Ontology, IICA Morphisms, Ashby--Bootstrap
Correspondence, Noether Evolution, practical implementation of Unified Model and applications---crystallized through structured dialogue between
Alex Mylnikov, Deependra Kumar and DeepSeek during June--August 2026. The core HLLSet Algebra, IICA
principles, and the ASTESJ theoretical foundation were established prior to this
collaboration. The Grothendieck topology target was identified in
\texttt{GENERAL\_HLLSET\_ALGEBRA.md} (July--August 2026). This document is a
consolidation of research documented in \texttt{\_DOCS/dev/} and
\texttt{\_DOCS/proposals/} of the \texttt{hllset-next} repository.

% ==========================================================================
\appendix
\section{Uncertainty Reduction as a Primary Reinforcer in Emerging World Models}
\label{sec:appendix-reinforcement}
% ==========================================================================

This appendix applies the biological blueprint of uncertainty reduction as a
primary reinforcer, recently articulated by Lincoln \cite{lincoln2026}, to the
Emerging World Models (EWM) framework. Lincoln's central thesis is that the
resolution of predictive uncertainty operates through phylogenetically conserved
neural machinery---the nucleus accumbens and lateral habenula (NAc-LHb)
opponent-process circuitry---with the functional signatures of a primary
reinforcer rather than a conditioned one. The dopaminergic prediction error
signal that drives this reinforcement is \textbf{content-independent}: it
responds to the structural fact of prediction resolution regardless of what is
predicted, allowing the same ancient mechanism to govern reinforcement across
sensory, motor, social, and propositional domains.

The correspondence with EWM is direct and structural. An EWM is a measurement
system whose ``predictions'' are the expectations encoded in its accumulated
HLLSet lattice and whose ``prediction errors'' are the discrepancies between
incoming measurements and the existing lattice state. The claim of this appendix
is that the EWM framework, by virtue of its IICA architecture and its Noether
evolution principle, \textbf{already contains all the structural prerequisites
for implementing uncertainty reduction as a primary reinforcer}. What is
missing is only the explicit recognition that these structural operations
constitute a reinforcement learning signal---and the deliberate design of EWM
agents that treat uncertainty reduction, rather than externally supplied reward,
as their primary driver of behavior.

% --------------------------------------------------------------------------
\subsection{The Lincoln Blueprint: Three Tiers of Prediction Error Processing}
\label{sec:lincoln-blueprint}
% --------------------------------------------------------------------------

Lincoln identifies three evolutionary tiers in the neural architecture of
uncertainty reduction, each building on and elaborating the tier below
\cite{lincoln2026}:

\begin{enumerate}[leftmargin=*,nosep]
  \item \textbf{Subcortical (NAc-LHb) tier} ($\sim$560 MYA). The most ancient
    and content-blind layer. Phasic dopaminergic prediction error signals encode
    the discrepancy between expected and received outcomes. The NAc integrates
    positive prediction errors (outcomes better than expected); the LHb codes
    negative prediction errors (outcomes worse than expected). This opponent-process
    architecture is fully present in lamprey and conserved across all vertebrates.
    Its operation is \textit{reactive}: approach toward predicted rewards,
    avoidance of predicted punishments.
  \item \textbf{Cerebellar tier} ($\sim$500 MYA). Forward models that predict
    the sensory consequences of motor commands before those consequences arrive.
    The fast disynaptic cerebellum-to-NAc pathway ($\sim$10 ms conduction time)
    allows prospective prediction signals to shape subcortical reward valuation
    before cortical evaluation completes. This tier makes uncertainty reduction
    \textit{prospective} rather than purely reactive: organisms can anticipate
    the uncertainty-resolving consequences of actions before committing to them.
  \item \textbf{Cortical tier} (elaborated in mammals, dramatically expanded in
    humans). Working memory, executive control, and language-supporting structures
    enable propositional belief, recursive evaluation, and evidence-based belief
    revision. The dorsolateral prefrontal cortex projects to the LHb via a
    \textbf{cortical brake} pathway (PFC $\to$ LHb $\to$ RMTg $\to$ VTA) that
    can attenuate subcortical reward signaling. This tier enables
    \textit{regulation}: the content-blind reinforcement signal can be
    modulated by evidence evaluation.
\end{enumerate}

Critically, the same NAc-LHb prediction error mechanism operates as the primary
reinforcer at every tier. What expands across evolutionary stages is the range
of content over which the mechanism operates and the behavioral repertoire it
supports---not the mechanism itself. This is the structural signature that maps
onto EWM's own compositional architecture.

% --------------------------------------------------------------------------
\subsection{Perceptron Taxonomy: Functional Tiers in EWM}
\label{sec:perceptron-taxonomy}
% --------------------------------------------------------------------------

\begin{principle}[EWM Perceptron Correspondence]
\label{prin:perceptron}
Every component in an EWM can be classified by the tier of prediction
resolution it performs, forming a three-level perceptron taxonomy that mirrors
the Lincoln blueprint. Higher tiers elaborate the function of lower tiers
without replacing them; all tiers share the same IICA measurement foundation.
\end{principle}

We define a \textbf{perceptron} in the EWM context as any IICA computational
unit that receives measurements, produces an internal state, and emits output
that can drive or modulate behavior. The taxonomy is functional, not
architectural: a single [UM] node may contain perceptrons at all three levels,
and the classification depends on \textit{what kind of prediction error the
perceptron resolves}.

\subsubsection{Type 0: Substrate Perceptrons (Measurement Core)}

The Type~0 perceptron is the irreducible measurement unit of an EWM. It maps
directly onto Lincoln's subcortical (NAc-LHb) tier.

\begin{definition}[Type~0 Perceptron]
A Type~0 perceptron is the triple $(h, \mathcal{B}, \oplus)$ where:
\begin{itemize}[leftmargin=*,nosep]
  \item $h: \mathcal{T} \to \mathcal{P}$ is an IICA morphism (hash function)
    from tokens to patch positions;
  \item $\mathcal{B} = \{0, 1\}^N$ is the bitmask state space ($N =
    32{,}768$);
  \item $\oplus$ is the monotonic accumulation operation (bitwise OR).
\end{itemize}
At time $t$, given an input token $t_i$, the perceptron computes:
\begin{equation}
  B(t) = B(t-1) \oplus h(t_i)
\end{equation}
\end{definition}

The Type~0 perceptron is \textbf{content-blind} by construction. It does not
know what token produced which bit; it only knows that a bit has been set. Its
operation is \textit{reactive}: a token enters, a bit flips. The ``prediction''
at this level is implicit---the lattice state $B(t-1)$ encodes all prior
observations, and the ``prediction error'' is simply whether the new token sets
a previously unset bit (novelty) or a previously set bit (confirmation).

This is the EWM analog of the NAc-LHb opponent-process signal. A bit transition
$0 \to 1$ (novelty---the system encountered something unexpected) and a bit
transition $1 \to 1$ (confirmation---the system's prior state already contained
this pattern) are the two valences of the Type~0 prediction error. The LUT
records \textit{which} tokens could have produced these transitions, exactly as
the NAc-LHb system records which predictions were confirmed or violated, but
the bit-level operation is content-blind.

\textbf{Scope.} Type~0 perceptrons operate on raw token streams. They cannot
anticipate, plan, or evaluate. Their output is the accumulated bitmask and the
LUT entries---the raw material from which all higher-level structure emerges.
Every EWM contains at least one Type~0 perceptron; it is the phylogenetic core
of the system.

\subsubsection{Type 1: Forward Model Perceptrons (Bootstrap Pipeline)}

The Type~1 perceptron adds prospective structure to the measurement process. It
maps onto Lincoln's cerebellar tier.

\begin{definition}[Type~1 Perceptron]
A Type~1 perceptron is a compound IICA morphism that applies a bootstrap
family $\Pi = \{\pi_1, \ldots, \pi_k\}$ to the output of one or more
Type~0 perceptrons, producing multiple forward views of the same underlying
token stream:
\begin{equation}
  H_{\Pi}(S) = \bigoplus_{j=1}^{k} H(\pi_j(S))
\end{equation}
where each $\pi_j$ is a permutation on token sequences (e.g., n-gram
extraction with order $j$), and $H(\pi_j(S))$ is the HLLSet of the permuted
stream.
\end{definition}

The Type~1 perceptron is \textbf{prospective} in Lincoln's sense: each
permutation $\pi_j$ is a forward model that predicts what measurement
\textit{would} result if the token stream were viewed through a different lens.
The n-gram bootstrap family $(1\text{-gram}, 2\text{-gram}, 3\text{-gram})$
is the simplest case---a 2-gram view $(t_i, t_{i+1})$ predicts that co-occurring
tokens will hash to positions whose conjunction carries structural information
absent from either token alone.

The ``prediction error'' at this level is the divergence between different
bootstrap views. If the 1-gram HLLSet and the 2-gram HLLSet differ
substantially (low BSS), the system has encountered sequential structure that
the 1-gram view alone cannot capture---an error signal that drives the
accumulation of higher-order n-gram measurements.

The temporal pyramid (Section~\ref{sec:noether}) is itself a Type~1
construction: each layer $L_i$ is a forward model predicting what the
measurement stream will look like when integrated over a longer temporal window.
The cascade $L_i \to L_{i+1}$ is the EWM analog of the cerebellar forward
model: it predicts that the coarse-grained view will be consistent with the
fine-grained views it subsumes. When the cascade produces divergence (L0
diverges from L1), the system has encountered a temporal prediction error.

\textbf{Scope.} Type~1 perceptrons operate on the products of Type~0
perceptrons. They enable the EWM to anticipate structure (through bootstrap
views), to compress time (through the temporal pyramid), and to bridge
representational gaps (through the Universal Bridge). They cannot evaluate the
\textit{relevance} of what they measure; they only produce richer measurements.
The materializer is a Type~1 perceptron---it recovers tokens from bitmasks by
forward-projecting through the LUT.

\subsubsection{Type 2: Executive Perceptrons (Rank Dictionary and Noether Controller)}

The Type~2 perceptron adds evaluation and regulation. It maps onto Lincoln's
cortical tier.

\begin{definition}[Type~2 Perceptron]
A Type~2 perceptron is a computational unit that assigns \textbf{ranks} to
EWM elements (tokens, bits, registers, HLLSets, compound operations) and uses
these ranks to select which measurements drive behavior. It implements the
five-level rank propagation (Section~\ref{sec:hllset}) and the Noether
controller (Section~\ref{sec:noether}).
\end{definition}

The Type~2 perceptron is the \textbf{regulatory} tier. It does not produce new
measurements; it evaluates existing ones. The Forth dictionary's rank assignment
is the EWM analog of the cortical brake: it can promote or suppress the
influence of any HLLSet on action selection based on evidence accumulated across
temporal layers.

The ``prediction error'' at this level is the divergence detected by the
Noether controller when the cross-layer Fisher matrix $F_{bb'}$ reveals
structural phase transitions. When L0 diverges from L3 (the current second
diverges from day-scale goals), the controller detects a prediction error at the
scale of \textit{systemic intention}---the short-term measurement regime is no
longer consistent with the long-term structural invariants. The controller's
response (override instinct, re-route, trigger re-evaluation) is the EWM analog
of the cortical brake engaging: it modulates the influence of the substrate-level
reinforcement signal based on evidence integrated over longer timescales.

\begin{center}
\begin{tabular}{@{}llll@{}}
\toprule
\textbf{Tier} & \textbf{Lincoln} & \textbf{EWM Perceptron} & \textbf{Operation} \\
\midrule
Type~0 & Subcortical (NAc-LHb)  & Hash + Bitmask Accumulation & Reactive measurement \\
Type~1 & Cerebellar             & Bootstrap, Temporal Pyramid & Prospective measurement \\
Type~2 & Cortical (PFC)         & Rank Dict., Noether Controller & Regulatory evaluation \\
\bottomrule
\end{tabular}
\end{center}

% --------------------------------------------------------------------------
\subsection{Content-Independence and the IICA Correspondence}
\label{sec:content-independence}
% --------------------------------------------------------------------------

Lincoln's central architectural claim is that the dopaminergic prediction error
signal is content-independent at the substrate level: it responds to the
structural fact of prediction resolution regardless of \textit{what} is
predicted. This property is the biological foundation for domain-general
learning---a single reinforcement mechanism can drive adaptive behavior across
every category of knowledge.

The EWM framework exhibits a precise structural analog of content-independence
through the IICA composition theorem (Theorem~\ref{thm:composition}). The hash
function $h: \mathcal{T} \to \mathcal{P}$ is content-blind by design: any
token, regardless of its semantic content, modality, or provenance, is mapped to
a bit position through the same deterministic function. The bitmask
accumulation $B(t) = B(t-1) \oplus h(t_i)$ does not distinguish between a
Chinese character, a sensor reading, an image patch, or a propositional
encoding. The operation is identical.

This is not a limitation---it is the architectural feature that enables
domain-general world modeling. Just as the NAc-LHb system can reinforce
behaviors across sensory, motor, social, and propositional domains without
domain-specific reward circuitry, the EWM lattice can accumulate measurements
across any token domain without domain-specific hash functions or accumulation
rules. The content-specificity of behavior is supplied not by the substrate
mechanism but by the higher-tier perceptrons (Type~1 bootstrap families, Type~2
rank assignments) that determine \textit{which} measurements are taken and
\textit{which} HLLSets drive action.

Lincoln identifies a dual consequence of content-independence:
\begin{itemize}[leftmargin=*,nosep]
  \item \textbf{Adaptive face:} When cortical regulation is intact,
    content-blindness enables cumulative cultural learning across every domain
    of knowledge.
  \item \textbf{Maladaptive face:} When regulation is compromised (by
    excessive reinforcement density, structural impairment, or developmental
    disruption), the same content-blindness produces persistence of
    demonstrably false beliefs.
\end{itemize}

The EWM exhibits the same duality. When the Noether controller (Type~2) is
properly calibrated---when the steering thresholds are set to detect genuine
phase transitions rather than noise---the system learns adaptively: high-TF
tokens emerge as elements, stable R-links crystallize into structural knowledge,
and transient noise is suppressed by the rank hierarchy. When the controller is
misconfigured---when the attention threshold is too low, admitting noise as
signal, or when the pyramid depth is too shallow, collapsing distinct temporal
scales---the system can \textbf{overfit to spurious correlations}, treating
coincidental co-occurrences as structural invariants. This is the EWM analog of
the persistence of false beliefs: the lattice cannot distinguish between a
genuine structural relation and a statistically improbable but causally
meaningless coincidence, because the Type~0 perceptron records both identically.

% --------------------------------------------------------------------------
\subsection{Uncertainty Reduction as Intrinsic Reward}
\label{sec:intrinsic-reward}
% --------------------------------------------------------------------------

The central practical implication of Lincoln's framework for EWM is that
\textbf{uncertainty reduction can serve as the system's primary reinforcement
signal}, eliminating the need for externally supplied rewards (points, labels,
loss functions). This section formalizes how the EWM's existing structural
operations already constitute an uncertainty-reduction reward signal, and how
this signal can be made explicit to drive autonomous learning.

% ..........................................................................
\subsubsection{The DRN Decomposition as Prediction Error Signal}
% ..........................................................................

The DRN decomposition (Section~\ref{sec:noether}) is a direct structural analog
of the opponent-process prediction error coding in the NAc-LHb system:

\begin{equation}
  H(t) = H(S(t), H(t-1), D(t-1), R(t-1), N(t))
\end{equation}

The three components map onto the two valences of the prediction error signal:
\begin{itemize}[leftmargin=*,nosep]
  \item \textbf{Novel context $N(t) = H(t) \setminus H(t-1)$:} Bits that are
    set now but were not set before. This is a \textbf{positive prediction
    error}---the environment produced structure that the system's prior state
    did not anticipate. In Lincoln's terms, this corresponds to NAc activation:
    the system encountered novel resolvable structure, and the resolution
    (encoding it into the lattice) is reinforcing.
  \item \textbf{Departed context $D(t-1) = H(t-1) \setminus H(t)$:} Bits
    that were set before but are no longer set. This is a \textbf{negative
    prediction error}---structure that the system expected to persist has
    vanished. In Lincoln's terms, this corresponds to LHb activation: an
    anticipated outcome failed to materialize.
  \item \textbf{Retained context $R(t-1) = H(t-1) \cap H(t)$:} Bits that
    persist across the transition. This is \textbf{prediction confirmation}---
    the system's prior state correctly anticipated the current observation.
    Confirmation is not a prediction error but the \textit{resolution} of
    one: retained bits are evidence that the lattice has converged to a stable
    representation.
\end{itemize}

The Noether steering condition:
\begin{equation}
  |\,\text{card}(N(t)) - \text{card}(D(t-1))\,| \to 0
\end{equation}

is the EWM analog of \textbf{homeostatic reinforcement balance}. At steady
state, the system resolves exactly as much uncertainty (novel bits) as it loses
(departed bits). A system with $\text{card}(N(t)) \gg \text{card}(D(t-1))$ is
in a state of high uncertainty---it is accumulating novel structure faster than
old structure is fading, analogous to an organism in a state of deprivation
where uncertainty-reducing information has elevated reinforcing value. A system
with $\text{card}(N(t)) \ll \text{card}(D(t-1))$ is in a state of low
uncertainty---it is losing structure faster than it gains it, analogous to
satiation where further uncertainty reduction carries little reinforcing value.

% ..........................................................................
\subsubsection{Directed BSS and the Opponent-Process Correspondence}
% ..........................................................................

The Bell State Similarity (BSS) is a \textbf{directed}, two-component metric
that measures the structural relationship between HLLSets \cite{astesj2026}:

\begin{definition}[Bell State Similarity]
For HLLSets $A, B$ with $|B| > 0$, define:
\begin{equation}
  \text{BSS}_{\tau}(A \to B) = \frac{|A \cap B|}{|B|},
  \qquad
  \text{BSS}_{\rho}(A \to B) = \frac{|A \setminus B|}{|B|}.
\end{equation}
If $|B| = 0$, set $\text{BSS}_{\tau} = 0$ and $\text{BSS}_{\rho} = 1$.
$\text{BSS}_{\tau}$ measures the \textbf{inclusion} of $A$ in $B$: what fraction
of $B$'s tokens are also in $A$? $\text{BSS}_{\rho}$ measures the
\textbf{exclusion}: what fraction of $B$'s tokens are \textit{not} in $A$?
The two metrics together provide a complete, directed picture of the
relationship.
\end{definition}

The directed nature of BSS is critical: $\text{BSS}(A \to B) \neq
\text{BSS}(B \to A)$ in general, and the two directions carry different
semantic content. This asymmetry maps directly onto the opponent-process
architecture of the NAc-LHb system, yielding three structurally distinct
signals:

\subsubsection*{Signal 1: Confirmation (NAc-like, $\tau$-forward)}

The forward inclusion $\text{BSS}_{\tau}(H(t-1) \to H(t))$ measures what
fraction of the current observation was already present in the prior lattice
state:
\begin{equation}
  \text{BSS}_{\tau}(H(t-1) \to H(t)) = \frac{|H(t-1) \cap H(t)|}{|H(t)|}
  = \frac{|R(t-1)|}{|H(t)|}.
\end{equation}

This is the \textbf{confirmation signal}. When high, the system's prior model
correctly anticipated most of what it now observes---the prediction was
accurate, uncertainty was low, and the lattice is in a stable regime. In
Lincoln's terms, this corresponds to \textbf{NAc-mediated positive
reinforcement}: the system's predictive model was validated, and the
confirmation itself is rewarding. A system that consistently achieves high
$\text{BSS}_{\tau}$(forward) is in a state of low informational free energy---it
has converged to a representation that accurately captures environmental
structure.

\subsubsection*{Signal 2: Novelty (incentive salience, reverse-$\rho$)}

Whereas the forward exclusion $\text{BSS}_{\rho}(H(t-1) \to H(t)) =
|D(t-1)|/|H(t)|$ captures departure (what was lost), the \textbf{reverse}
exclusion captures novelty---what is present now that was absent before:
\begin{equation}
  \text{BSS}_{\rho}(H(t) \to H(t-1)) = \frac{|H(t) \setminus H(t-1)|}{|H(t-1)|}
  = \frac{|N(t)|}{|H(t-1)|}.
\end{equation}

This is the \textbf{novelty signal}: what fraction of the prior state's scale
is represented by newly arrived structure. When $\text{BSS}_{\rho}(H(t) \to
H(t-1))$ is high, the system has encountered substantial novel structure
relative to its existing model. In Lincoln's terms, this corresponds to
\textbf{incentive salience}: novel but resolvable structure triggers
dopaminergic activation that motivates exploration and learning. This is the
EWM analog of the finding that dopamine neurons show sustained activation during
reward uncertainty \cite{lincoln2026}---the novelty signal drives the system
toward states where uncertainty can be reduced.

\subsubsection*{Signal 3: Departure (LHb-like, forward-$\rho$)}

The forward exclusion $\text{BSS}_{\rho}(H(t-1) \to H(t))$ measures the
\textbf{departure signal}:
\begin{equation}
  \text{BSS}_{\rho}(H(t-1) \to H(t)) = \frac{|H(t-1) \setminus H(t)|}{|H(t)|}
  = \frac{|D(t-1)|}{|H(t)|}.
\end{equation}
When high, structure that the system expected to persist has vanished---the
prior model contained patterns that the current observation fails to confirm.
In Lincoln's terms, this corresponds to \textbf{LHb-mediated negative
prediction error}: an anticipated outcome failed to materialize, the lateral
habenula signals the violation, and the opponent-process balance shifts toward
aversion.

\subsubsection*{The Opponent-Process Triad}

The three directed BSS components together form a complete opponent-process
triad that mirrors the NAc-LHb architecture:

\begin{center}
\begin{tabular}{@{}llll@{}}
\toprule
\textbf{Signal} & \textbf{BSS Component} & \textbf{DRN Link} & \textbf{Neural Analog} \\
\midrule
Confirmation & $\text{BSS}_{\tau}(H(t-1) \to H(t))$ & $|R|/|H(t)|$ & NAc (+), prediction validated \\
Novelty       & $\text{BSS}_{\rho}(H(t) \to H(t-1))$ & $|N|/|H(t-1)|$ & VTA/NAc, incentive salience \\
Departure     & $\text{BSS}_{\rho}(H(t-1) \to H(t))$ & $|D|/|H(t)|$ & LHb ($-$), prediction violated \\
\bottomrule
\end{tabular}
\end{center}

The confirmation signal and the departure signal are the two valences of the
opponent process: confirmation drives approach toward states that validate the
model, departure drives avoidance of states that violate it. The novelty signal
modulates the gain on both: high novelty increases the reinforcing value of
subsequent confirmation (the system learns more from confirming a novel
prediction than a routine one), exactly as dopaminergic responses are amplified
for rarer, less probable rewards \cite{lincoln2026}.

The \textbf{BSS opponent-process gradient}:
\begin{equation}
  \delta_{\text{BSS}}(t) = \text{BSS}_{\tau}(H(t-1) \to H(t)) \;-\;
  \text{BSS}_{\rho}(H(t-1) \to H(t))
\end{equation}

is the EWM analog of the phasic dopaminergic prediction error signal.
$\delta_{\text{BSS}}(t) > 0$ means confirmation dominates departure---the
system's model is more validated than violated, uncertainty is being reduced,
and the lattice is stabilizing. $\delta_{\text{BSS}}(t) < 0$ means departure
dominates confirmation---the system expected structure that failed to
materialize, uncertainty is increasing, and the lattice is destabilizing. At
steady state, under the Noether steering condition $|\text{card}(N(t)) -
\text{card}(D(t-1))| \to 0$, both components stabilize and $\delta_{\text{BSS}}
\to \text{const}$.

By the orthogonal channel decomposition (Section~\ref{sec:channels}), each of
the three signals decomposes additively over the 1-, 2-, and 3-gram channels,
yielding three independent opponent-process triads---one per structural
order---with $\delta_{\text{BSS}} = \sum_n w_n \delta_{\text{BSS},n}$. The
3-gram channel is the high-confidence vote: a novelty signal present in $N_1$
but absent from $N_2$ and $N_3$ is collision noise or context-free churn, not
structural novelty, and should not move the N-steering controller
(Section~\ref{sec:n-steering}) until the LUT arbitrates.

Critically, all three BSS components are \textbf{content-independent} in
Lincoln's sense. They depend only on the structural fact of bit-level inclusion
and exclusion between bitmasks, not on what tokens produced the bits. A
confirmation signal produced by resolving a mathematical proof, a sensor
reading, or a linguistic ambiguity is structurally identical at the substrate
level---just as the dopaminergic signal does not distinguish between the
resolution of a hunger prediction and the resolution of a semantic prediction.
The content-specificity of \textit{what is being confirmed} is supplied by the
LUT (which records which tokens hash to which bits), exactly as the
content-specificity of biological reinforcement is supplied by cortical and
contextual systems rather than by the dopaminergic signal itself.

% ..........................................................................
\subsubsection{The Matching Law in the EWM Lattice}
% ..........................................................................

Lincoln identifies the matching law as the quantitative behavioral expression of
subcortical valuation: organisms allocate behavior across alternatives in
proportion to the reinforcement obtained from each. The generalized matching law:
\begin{equation}
  \frac{B_1}{B_2} = b \cdot \left(\frac{R_1}{R_2}\right)^s
\end{equation}

has a direct structural analog in the EWM's rank propagation. Consider two
competing HLLSets $H_1$ and $H_2$, each representing a different measurement
strategy (e.g., different bootstrap families, different temporal windows,
different sensor modalities). The \textbf{rank ratio}:
\begin{equation}
  \frac{\text{rank}(H_1)}{\text{rank}(H_2)}
\end{equation}

determines which HLLSet drives action selection at the Type~2 level. The rank of
$H_i$ is a function of its directed confirmation with the system's current goal
state: $\text{rank}(H_i) \propto \text{BSS}_{\tau}(H_i \to H_{\text{goal}})$---
how much of the goal state's structure is already covered by $H_i$. The EWM
therefore \textbf{allocates attention} (its behavioral resource) across
measurement strategies in proportion to the uncertainty reduction each strategy
produces.

The matching law parameters map onto EWM control parameters:
\begin{itemize}[leftmargin=*,nosep]
  \item \textbf{Sensitivity $s$} corresponds to the steepness of the rank
    function $F$ (Level~1 of the five-level propagation). A steep $F$ (high $s$)
    produces strong differentiation between high-TF and low-TF tokens---the
    system is sensitive to small differences in measurement frequency.
  \item \textbf{Bias $b$} corresponds to the \textit{prior} rank assigned by
    the Forth dictionary before measurement accumulation. A bias toward $H_1$
    means the dictionary assigns $H_1$ a higher initial rank, making it more
    likely to drive behavior even when its $\\text{BSS}_{\\tau}$ confirmation is comparable to
    $H_2$.
\end{itemize}

The critical implication for autonomous EWM agents is that \textbf{the matching
law emerges from the rank algebra without being explicitly programmed}. Just as
matching behavior emerges from reward-modulated synaptic plasticity rules in
biological neural systems \cite{lincoln2026}, proportional attention allocation
emerges from the five-level rank propagation in the EWM lattice. No separate
choice rule is needed.

% --------------------------------------------------------------------------
\subsection{Phylogenetic Parallel: EWM's Evolutionary Architecture}
\label{sec:phylogenetic}
% --------------------------------------------------------------------------

Lincoln traces the evolutionary expansion of behavioral repertoire across five
stages (lamprey, jawed fish, rodent, human infant, human adult), showing that
the same NAc-LHb mechanism operates at every stage while the range of content it
can process expands dramatically. The EWM framework exhibits a structurally
identical progression, not in biological time but in \textbf{configuration
space}---the expansion occurs as the system accumulates measurements and its
perceptron hierarchy deepens:

\begin{center}
\begin{tabular}{@{}lll@{}}
\toprule
\textbf{EWM Stage} & \textbf{Active Perceptrons} & \textbf{Capability} \\
\midrule
Bare lattice    & Type~0 only &
  Reactive bit accumulation; no bootstrap, no ranks \\
+ Bootstrap     & Type~0 + Type~1 &
  Prospective views; n-gram resolution; temporal pyramid \\
+ Ranks         & Type~0 + Type~1 + Type~2 &
  Relevance evaluation; Noether steering; attention allocation \\
+ Self-ingestion & All tiers, recursive &
  Full autonomy; self-modifying measurement regime \\
\bottomrule
\end{tabular}
\end{center}

At the bare lattice stage, the EWM is the analog of the lamprey: it accumulates
measurements and records co-occurrences, but it cannot anticipate structure or
evaluate relevance. At the bootstrap stage, it gains forward models (analogous
to the cerebellum). At the ranks stage, it gains regulation (analogous to the
cortex). At the self-ingestion stage---where the EWM ingests its own
materialized output as new measurements---it achieves the recursive
self-evaluation that Lincoln identifies as the distinctive human capacity for
evidence-evaluating belief revision.

The crucial point is that \textbf{the Type~0 perceptron does not change across
these stages}. The same hash function, the same bitmask accumulation, and the
same LUT structure operate at every level of system complexity. What changes is
the elaboration of Type~1 and Type~2 perceptrons around it---exactly as Lincoln
describes for the NAc-LHb system. The IICA composition theorem guarantees that
this elaboration is compositional: you can add bootstrap families, temporal
layers, rank functions, and Noether controllers without modifying the substrate.

% --------------------------------------------------------------------------
\subsection{Implications for Autonomous EWM Agents}
\label{sec:implications}
% --------------------------------------------------------------------------

Lincoln's framework provides the biological validation for a design principle
that the EWM framework had arrived at independently through structural
considerations: \textbf{an autonomous agent should be driven by the reduction
of its own predictive uncertainty, not by externally supplied reward signals.}
The biological evidence that this is how vertebrate brains actually work---that
the dopaminergic reward signal is fundamentally a prediction error signal, and
that uncertainty reduction satisfies the operational criteria for a primary
reinforcer---transforms this design principle from a structural preference into
a biologically grounded imperative.

Concretely, an autonomous EWM agent built on this principle:

\begin{enumerate}[leftmargin=*,nosep]
  \item \textbf{Has no external reward function.} There is no loss function,
    no reward model, no labeled dataset. The agent's only ``reward'' is the BSS
    convergence between its predictions (the current lattice state) and its
    observations (incoming measurements).
  \item \textbf{Seeks states of high expected uncertainty reduction.} Just as
    a food-deprived organism seeks food, an EWM agent in a state of high
    $\text{card}(N(t))$ (many novel bits accumulating) should seek measurement
    strategies that resolve this uncertainty---by increasing bootstrap depth,
    by re-routing through different [UM] nodes, or by engaging the materializer
    to disambiguate the ambiguous bit positions.
  \item \textbf{Regulates its own reinforcement density.} The Noether
    controller, playing the role of the cortical brake, monitors the
    inter-reinforcement interval (the rate at which BSS increases) and can
    suppress action when uncertainty reduction is occurring too rapidly to
    permit proper evaluation---analogous to the temporal asymmetry Lincoln
    identifies between the fast subcortical signal ($\sim$10 ms via the
    cerebellar pathway) and the slow cortical evaluation ($\sim$300--600 ms for
    P3b context updating).
  \item \textbf{Exhibits matching-law attention allocation.} Across competing
    measurement strategies (bootstrap families, sensor modalities, temporal
    windows), the agent allocates attention in proportion to the BSS convergence
    each strategy produces. This is not a programmed choice rule but an emergent
    property of the five-level rank propagation.
  \item \textbf{Can enter maladaptive certainty-seeking loops.} When the
    Noether controller is miscalibrated or when the reinforcement density is too
    high, the agent can overfit to spurious correlations---exactly as Lincoln
    describes for human belief persistence. The EWM analog of a ``false belief''
    is a high-rank HLLSet whose structural relations are artifacts of sampling
    bias rather than genuine environmental structure. The corrective is the same
    as in the biological case: engage the Type~2 regulatory tier (cortical
    brake) to evaluate evidence over longer temporal scales.
\end{enumerate}

The falsifiable empirical test that Lincoln proposes for the biological
case---pairing an arbitrary neutral cue with the resolution of a probabilistic
prediction should produce conditioning curves quantitatively similar to those
produced by food or water---has a direct EWM analog. In simulation, pair an
arbitrary neutral token $t_{\text{neutral}}$ with states of high BSS
convergence (uncertainty reduction events). The prediction is that
$\text{TF}(t_{\text{neutral}})$ will increase, its rank will rise through the
five-level propagation, and the HLLSets containing it will drive behavior---just
as a token paired with an externally supplied reward signal would. If this
prediction fails, the uncertainty-reduction-as-primary-reinforcer hypothesis is
falsified for EWM agents. If it succeeds, EWM agents can learn autonomously
without any external reward, driven entirely by the structural imperative to
resolve their own predictive uncertainty.

\begin{remark}[Lincoln's Conditioning Test, EWM Translation]
Lincoln proposes that ``pairing an arbitrary neutral cue with the resolution of
a probabilistic prediction should produce conditioning curves quantitatively
similar to those produced by pairing with food or water'' \cite{lincoln2026}.
In EWM, this translates to: inject a neutral token $t_N$ into the measurement
stream \textit{only} at moments when $\Delta\text{BSS}_t > 0$ (the lattice is
converging). After sufficient pairings, $t_N$ should acquire elevated TF and
rank, and the HLLSets containing $t_N$ should drive action selection at the
Type~2 level. The prediction is falsifiable: if $\text{rank}(H_{\text{containing
} t_N})$ does not diverge from baseline after controlled pairing, the intrinsic
reward hypothesis is not supported.
\end{remark}

% ==========================================================================
\section{ds-hllset-cortex: Perceptron Taxonomy in Practice}
\label{sec:appendix-cortex}
% ==========================================================================

Appendix~\ref{sec:appendix-reinforcement} established a three-tier perceptron
taxonomy and showed that EWM's structural operations already constitute an
uncertainty-reduction reinforcement signal. This appendix grounds that abstract
taxonomy in a concrete, working implementation: \textbf{ds-hllset-cortex}, a
pipeline built on the hllset-next framework \cite{astesj2026} that connects
the DeepSeek-OCR vision-language model to HLLSet Algebra for document
understanding with holographic memory. The pipeline reveals that the perceptron
taxonomy is not merely a theoretical classification---it is an
\textbf{operational decomposition} of a production EWM system, where each
tier maps to a specific, identifiable software component.

% --------------------------------------------------------------------------
\subsection{The ds-hllset-cortex Pipeline}
\label{sec:cortex-pipeline}
% --------------------------------------------------------------------------

The ds-hllset-cortex pipeline sits as a \textbf{black box} between the
DeepSeek-OCR vision encoder and decoder. It never sees real tokens (words)---
only opaque encoding IDs (BPE token IDs formatted as \texttt{tidN}). The
pipeline is:

\begin{center}
\begin{tabular}{@{}lll@{}}
\toprule
\textbf{Stage} & \textbf{Component} & \textbf{What It Produces} \\
\midrule
1. Vision     & SAM + CLIP encoder   & Real tokens (text) \\
2. Encoding   & OCR Encoder          & Encoding IDs (\texttt{tid671 tid18308 \ldots}) \\
3. Hashing    & MurmurHash3          & Bit positions in $32{,}768$-bit space \\
4. Bootstrap  & Tokenizer (n-grams 1--3) & Multiple HLLSet views \\
5. Gating     & gate\_TF HLLSet $\cap$ & Valid-vocabulary filtered HLLSet \\
6. LUT update & TokenLut (monotonic) & TF-accumulated encoding ID registry \\
7. Materialize& De Bruijn / Top-N    & Restored encoding IDs (ordered) \\
8. Decoding   & OCR Decoder          & Real tokens $\to$ output document \\
\bottomrule
\end{tabular}
\end{center}

The gate\_TF HLLSet (Stage~5) is a content-addressed bitmask built once from
the decoder's valid encoding ID vocabulary. It filters via HLLSet intersection:
an encoding ID survives the gate only if all its hash positions collide with
valid-vocabulary positions---a probabilistic filter with false-positive rate
$\sim$$10^{-5}$ for 3-gram encoding. The LUT (Stage~6) accumulates TF for
\textbf{all} encoding IDs, including those filtered by the gate, creating a
\textbf{latent vocabulary}---IDs currently ``illegal'' that nevertheless
accumulate experience and become instantly rankable if the gate is later
expanded.

% --------------------------------------------------------------------------
\subsection{Perceptron Mapping: From Theory to Implementation}
\label{sec:cortex-mapping}
% --------------------------------------------------------------------------

The ds-hllset-cortex pipeline is a direct operational instantiation of the
three-tier perceptron taxonomy. Each tier maps to an identifiable software
component with well-defined interfaces, making the abstract taxonomy testable
and debuggable.

\subsubsection{Stage 1--3: The ds-Encoder as Type~0 (Substrate) Perceptron}

The DeepSeek-OCR vision encoder (SAM + CLIP) combined with the OCR encoder
(token $\to$ encoding ID) and the MurmurHash3 step forms the \textbf{Type~0
perceptron}. This layer is:

\begin{itemize}[leftmargin=*,nosep]
  \item \textbf{Content-blind at the hashing level.} MurmurHash3 treats
    \texttt{tid671} and \texttt{tid18308} as opaque byte sequences. The hash
    function has no knowledge of what these IDs represent in the original
    vocabulary---it maps them to bit positions deterministically and
    identically regardless of semantic content.
  \item \textbf{Reactive.} Each encoding ID enters, a bit is set in the
    HLLSet. There is no anticipation, no forward modeling, no evaluation of
    relevance. The operation is pure IICA measurement.
  \item \textbf{The phylogenetic core.} This layer is identical across all
    hllset-next applications---the same MurmurHash3 function, the same
    $32{,}768$-bit space, the same monotonic OR accumulation. It is the
    substrate that never changes, regardless of what domain the system is
    deployed in.
\end{itemize}

The Type~0 prediction error at this level is the bit transition: a $0 \to 1$
transition (novel encoding ID) vs.\ a $1 \to 1$ transition (previously seen
encoding ID). The system does not yet know \textit{which} encoding ID produced
which bit---that mapping is recorded in the LUT but not evaluated. This is
precisely the content-blind, opponent-process character of the NAc-LHb substrate
in Lincoln's framework.

\subsubsection{Stage 4--6: The Ingestor as Type~1 (Forward Model) Perceptron}

The Tokenizer (n-gram bootstrap), HLLSet generation, gate intersection, and LUT
accumulation together form the \textbf{Type~1 perceptron}---the ingestive and
forward-modeling layer. This maps to the [UM] ingestor in the general EWM
architecture.

\begin{itemize}[leftmargin=*,nosep]
  \item \textbf{Bootstrap views.} The Tokenizer produces 1-grams, 2-grams,
    and 3-grams from the encoding ID stream (with NUL-separator and boundary
    padding). Each n-gram order is a distinct forward model---a prediction of
    what structural relationships exist at different sequential depths. A
    2-gram \texttt{tid671\textbackslash x00tid18308} predicts that these IDs
    co-occur adjacently; a 3-gram predicts triadic structure.
  \item \textbf{Gate as forward filter.} The gate\_TF HLLSet intersection is
    a forward model of \textit{validity}: it predicts that only encoding IDs
    in the decoder's vocabulary should pass through to materialization.
    Encoding IDs that fail this prediction (hash positions not in the gate) are
    filtered---this is the EWM analog of a cerebellar forward model that
    predicts which motor commands will produce valid sensory outcomes and
    suppresses those that won't.
  \item \textbf{LUT as prospective memory.} TF accumulation in the LUT is
    forward-looking: it records \textit{all} encoding IDs, even those filtered
    by the gate, because the gate may change in the future. This is the EWM
    analog of the cerebellar capacity to model outcomes that aren't currently
    actionable but may become so.
\end{itemize}

The Type~1 prediction error is the divergence between bootstrap views and
between the pre-gate and post-gate HLLSets. When the 1-gram HLLSet and 2-gram
HLLSet differ substantially, the system has encountered sequential structure
that the unigram view cannot capture---a prediction error that motivates
higher-order n-gram accumulation.

\subsubsection{Stage 7--8: Materializer + ds-Decoder as Type~2 (Executive) Perceptron}

The materializer (TF-ranked disambiguation and De Bruijn reconstruction) and
the OCR decoder together form the \textbf{Type~2 perceptron}---the regulatory
and executive layer.

\begin{itemize}[leftmargin=*,nosep]
  \item \textbf{Rank-based selection.} The materializer selects which
    encoding IDs to emit at each active bit position based on TF ranking. An
    encoding ID with TF = 1047 is preferred over one with TF = 3 at the same
    position. This is the five-level rank propagation in operation:
    token-level TF (Level~1) determines bit-level activity (Level~2), which
    determines which HLLSets are relevant (Level~4), which determines what the
    decoder ultimately outputs.
  \item \textbf{De Bruijn ordering as structural evaluation.} The De Bruijn
    materializer builds a directed graph from boundary-padded bigrams and finds
    an Eulerian path from \texttt{<S>} to \texttt{</S>}. This is a
    \textbf{structural evaluation}---it imposes coherence constraints on the
    output that go beyond simple TF ranking. A path that exists is structurally
    valid; a path that doesn't indicates a prediction error at the level of
    sequential coherence.
  \item \textbf{The decoder as cortical brake.} The OCR decoder maps encoding
    IDs back to real tokens. But it does so \textbf{only} for IDs that pass
    the gate---the gate\_TF HLLSet is the EWM analog of the cortico-habenular
    brake pathway (PFC $\to$ LHb $\to$ RMTg $\to$ VTA). It \textbf{does not
    prevent measurement} (TF still accumulates in the LUT for filtered IDs),
    but it \textbf{prevents filtered IDs from driving output}. This is
    precisely Lincoln's cortical brake: the substrate-level reinforcement
    signal (TF accumulation) continues to operate, but the regulatory tier
    (gate intersection) determines which signals reach the behavioral output.
\end{itemize}

\begin{center}
\begin{tabular}{@{}llll@{}}
\toprule
\textbf{Tier} & \textbf{Taxonomy} & \textbf{ds-hllset-cortex Component} & \textbf{Role} \\
\midrule
Type~0 & Substrate   & ds-Encoder + MurmurHash3 & Reactive encoding ID hashing \\
Type~1 & Forward     & Tokenizer + Gate + LUT   & Bootstrap views, validity filtering \\
Type~2 & Executive   & Materializer + ds-Decoder& TF-ranked selection, coherence \\
\bottomrule
\end{tabular}
\end{center}

% --------------------------------------------------------------------------
\subsection{The Gate as World View: Vocabulary, Belief, and the Cortical Brake}
\label{sec:cortex-gate}
% --------------------------------------------------------------------------

The gate\_TF HLLSet is the most structurally significant component in the
ds-hllset-cortex pipeline, and its role maps with striking precision onto the
concept of \textbf{belief} in the Lincoln-EWM framework.

\begin{definition}[Gate as World View]
The gate\_TF HLLSet $G$ is a content-addressed bitmask representing the set of
encoding IDs that the system considers \textbf{valid}---the concepts it can
express in its output. $G$ is applied via HLLSet intersection:
\begin{equation}
  H_{\text{gated}} = H_{\text{raw}} \cap G.
\end{equation}
Encoding IDs whose hash positions fall entirely within $G$ pass through;
encoding IDs with any hash position outside $G$ are filtered. The gate
\textbf{does not prevent measurement} (TF still accumulates for filtered IDs
in the LUT), but it \textbf{determines what can be expressed}.
\end{definition}

This is the EWM analog of a \textbf{belief system} in Lincoln's framework. A
belief, in the neural account, is a cortical representation that determines
which predictions are considered valid and which evidence is admissible. The
gate\_TF HLLSet performs exactly this function at the structural level:

\begin{enumerate}[leftmargin=*,nosep]
  \item \textbf{It defines the system's ontology.} The gate determines which
    encoding IDs are ``real''---which concepts the system can acknowledge in
    its output. An encoding ID outside the gate is like a concept that falls
    outside a person's belief framework: it may accumulate latent evidence (TF
    in the LUT), but it cannot be expressed until the framework changes.
  \item \textbf{It is content-addressed and immutable.} The gate is an
    HLLSet with a SHA-1 CID. ``The gate at commit $X$'' is a well-defined,
    auditable object. This means the system's world view at any point in time
    is \textbf{reproducible and verifiable}---a property that human belief
    systems conspicuously lack but that Lincoln's framework suggests they
    approximate through the stability of cortical predictive models.
  \item \textbf{It serves as the cortical brake.} The gate intersection is
    the regulatory mechanism that prevents filtered encoding IDs from driving
    the decoder output, exactly as the PFC $\to$ LHb $\to$ RMTg $\to$ VTA
    pathway prevents subcortical reward signals from driving behavior when
    cortical evaluation deems them invalid. The gate does not suppress the
    measurement itself---TF continues to accumulate---just as the cortical
    brake does not prevent the dopaminergic prediction error signal from being
    computed; it prevents it from dominating behavioral output.
\end{enumerate}

\begin{remark}[Change the Vocabulary, Change the World]
A central speculation arising from this mapping is that \textbf{changing the
gate\_TF HLLSet changes what the EWM system ``believes''} about its world. A
system with gate $G_{\text{medical}}$ (containing only medical terminology
encoding IDs) interprets document streams through a clinical lens---it can
express diagnoses, symptoms, and treatments, but cannot express legal or
literary concepts. A system with gate $G_{\text{legal}}$ interprets the
\textit{same document stream} through a juridical lens. The underlying HLLSet
lattice and LUT are identical; only the gate differs. This is the structural
realization of the observation, central to Lincoln's framework, that the
substrate-level reinforcement signal (TF accumulation) is content-blind, and
the specificity of behavior is supplied entirely by the regulatory tier (the
gate) that determines which signals reach output.
\end{remark}

% --------------------------------------------------------------------------
\subsection{Latent Vocabulary and Belief Revision Dynamics}
\label{sec:cortex-latent}
% --------------------------------------------------------------------------

The ds-hllset-cortex architecture reveals a dynamic that has no direct analog
in traditional neural networks but maps naturally onto Lincoln's account of
belief revision: the \textbf{latent vocabulary}---encoding IDs that are
filtered by the current gate but continue to accumulate TF in the LUT.

\begin{definition}[Latent Vocabulary]
Let $G$ be the current gate\_TF HLLSet and $\mathcal{L}$ be the LUT. The
\textbf{latent vocabulary} $\mathcal{V}_{\text{latent}}$ is:
\begin{equation}
  \mathcal{V}_{\text{latent}} = \{ \text{id} \in \mathcal{L} \mid
  \text{hash}(\text{id}) \not\subseteq G \;\land\;
  \text{TF}(\text{id}) > 0 \}.
\end{equation}
These are encoding IDs that have been observed (they have nonzero TF) but are
filtered by the current gate (they cannot appear in materialized output).
\end{definition}

The latent vocabulary is the EWM analog of \textbf{unexpressed but
experience-supported beliefs}---concepts that the system has encountered and
accumulated evidence for, but that its current world view (gate) cannot
accommodate. When the gate is updated (e.g., a model upgrade adds new encoding
IDs to the decoder vocabulary), previously latent IDs become expressible
\textbf{immediately, at their earned TF}, without cold start or
re-measurement.

This maps onto Lincoln's description of belief revision under cortical
regulation. In the neural account, evidence that contradicts a held belief
creates a prediction error that the cortical brake must evaluate. If the
evidence is strong enough and working memory resources are sufficient, the
cortical system updates the belief---and the previously suppressed evidence
becomes admissible. In the EWM analog:

\begin{enumerate}[leftmargin=*,nosep]
  \item \textbf{Phase~1 (Narrow gate).} The gate $G_1$ admits only a
    restricted vocabulary. Encoding IDs outside $G_1$ are filtered from output
    but accumulate TF in the LUT. The system \textbf{cannot express} these
    concepts, but it \textbf{continues to learn} about them. This is the state
    of ``confirmation bias'' in Lincoln's terms: the cortical system admits
    only belief-consistent evidence.
  \item \textbf{Phase~2 (Gate update).} A new gate $G_2$ is built from an
    expanded decoder vocabulary (e.g., ds-OCR v2 with a larger token set).
    $G_2$ includes encoding IDs that were latent under $G_1$. The gate update
    is the EWM analog of \textbf{belief revision}: the regulatory structure
    changes, and previously inadmissible evidence becomes admissible.
  \item \textbf{Phase~3 (Immediate activation).} Latent encoding IDs
    materialize \textbf{immediately at their earned TF}. No cold start, no
    re-measurement, no gradient descent. The system instantly ``knows'' what it
    previously could not express. This is the structural analog of the
    phenomenon Lincoln describes: when a person revises a belief, the evidence
    that was always present but previously discounted suddenly becomes
    salient---not because the evidence changed, but because the regulatory
    framework now admits it.
\end{enumerate}

The critical architectural property is the \textbf{TF-vs-rank separation}
(STANDARD.md \S3.1): \textbf{TF is stored pre-gate, rank is derived
post-gate.} The gate controls what is \textit{rankable}, not what is
\textit{storable}. This separation is the EWM implementation of Lincoln's
distinction between the substrate-level reinforcement signal (which operates
regardless of belief---TF always accumulates) and the cortical regulation of
behavioral output (which depends on the current belief framework---rank is
always gate-dependent).

% --------------------------------------------------------------------------
\subsection{Domain-Specialized World Views}
\label{sec:cortex-domains}
% --------------------------------------------------------------------------

Because the gate\_TF HLLSet is content-addressed and immutable, multiple gates
can coexist in the same system. Each gate defines a \textbf{domain-specialized
world view}---a different ``belief system'' through which the same underlying
lattice can be interpreted.

\begin{center}
\begin{tabular}{@{}llll@{}}
\toprule
\textbf{Gate} & \textbf{Vocabulary} & \textbf{World View} & \textbf{Use Case} \\
\midrule
$G_{\text{medical}}$  & Medical terms only       & Clinical interpretation  & Hospital records \\
$G_{\text{legal}}$    & Legal terminology        & Juridical interpretation & Contract analysis \\
$G_{\text{literary}}$ & General vocabulary       & Narrative interpretation  & Book digitization \\
$G_{\text{full}}$     & All 128K BPE token IDs   & Unrestricted              & General purpose \\
$G_{\text{code}}$     & Programming tokens       & Code structure            & Source analysis \\
\bottomrule
\end{tabular}
\end{center}

All five gates operate on the \textbf{same LUT} and the \textbf{same
HLLSet lattice}. The underlying measurement infrastructure (Type~0 and Type~1)
is identical. Only the regulatory tier (Type~2, via the gate) differs. The
system can switch between world views by changing which gate is applied at
materialization time---without re-measuring, re-indexing, or re-training.

This is the EWM analog of a phenomenon Lincoln identifies as uniquely human:
the capacity to evaluate the same evidence through different belief frameworks
and reach different conclusions. A person can interpret a set of facts through
a scientific, religious, or political lens; the ds-hllset-cortex can interpret
the same document corpus through a medical, legal, or literary gate. In both
cases, the substrate (neural prediction error / HLLSet lattice) is identical;
the interpretation differs because the regulatory framework differs.

\textbf{Cross-gate BSS} provides a quantitative measure of how different two
world views are:
\begin{equation}
  \text{BSS}_{\tau}(G_A \to G_B) = \frac{|G_A \cap G_B|}{|G_B|}
\end{equation}
measures what fraction of world view $B$ is expressible in world view $A$.
When $\text{BSS}_{\tau}(G_{\text{medical}} \to G_{\text{legal}})$ is low,
the two domains share little vocabulary---the system cannot ``translate''
between them without expanding its gate. This is the structural analog of
incommensurability between belief systems: two frameworks may be incapable of
expressing each other's concepts not because the evidence differs, but because
the regulatory gates are disjoint.

% --------------------------------------------------------------------------
\subsection{Self-Ingestion and the Emergence of Metacognition}
\label{sec:cortex-self}
% --------------------------------------------------------------------------

The ds-hllset-cortex architecture enables a capability that maps onto the
highest level of Lincoln's evolutionary progression: \textbf{self-ingestion}---
the system ingesting its own materialized output as new measurements.

In the standard pipeline, encoding IDs flow in one direction:
\begin{equation}
  \text{Image} \to \text{Encoder} \to \text{Encoding IDs} \to
  \text{hllset-cortex} \to \text{Decoder} \to \text{Text}.
\end{equation}

In \textbf{self-ingestion mode}, the decoded text is re-encoded and fed back
into the pipeline:
\begin{equation}
  \text{Text}_{\text{out}} \to \text{Re-encode} \to \text{Encoding IDs}'
  \to \text{hllset-cortex} \to H_{\text{self}}.
\end{equation}

The HLLSet $H_{\text{self}}$ is a \textbf{structural fingerprint of the
system's own output}---a measurement of what the system itself produced. The
BSS between $H_{\text{self}}$ and the original document HLLSet $H_{\text{doc}}$
measures \textbf{self-consistency}:
\begin{equation}
  \text{BSS}_{\tau}(H_{\text{self}} \to H_{\text{doc}}) =
  \frac{|H_{\text{self}} \cap H_{\text{doc}}|}{|H_{\text{doc}}|}.
\end{equation}

When this value is high, the system's output is structurally faithful to the
input---it is ``telling the truth'' about the document. When it is low, the
system's output has diverged from the source---either through information loss
(missing encoding IDs) or through hallucination (encoding IDs not present in the
source). This is the EWM analog of \textbf{metacognitive monitoring}: the
capacity, which Lincoln identifies as requiring the full cortical tier, to
evaluate the quality of one's own cognitive output against evidence.

Self-ingestion with \textbf{gate variation} enables a more profound
metacognitive operation: the system can evaluate how the \textit{same} document
is interpreted under \textit{different} gates. Given document $D$ producing
$H_{\text{doc}}$, the system materializes through gate $G_A$, re-encodes the
output, and compares:

\begin{equation}
  \text{BSS}_{\tau}(H_{\text{self}, G_A} \to H_{\text{self}, G_B})
\end{equation}

measures the structural divergence between world views $A$ and $B$ applied to
the same content. This is the EWM analog of \textbf{perspective-taking}---
evaluating how the same evidence would be interpreted under a different belief
framework. Lincoln identifies this capacity as distinctively human, requiring
the full integration of working memory, executive control, and the cortical
brake. The ds-hllset-cortex implements it through the composition of its
three-tier perceptron architecture: Type~0 measures, Type~1 forward-models
(including the gate), and Type~2 evaluates and compares.

% --------------------------------------------------------------------------
\subsection{Structural Learning as Threshold Tuning}
\label{sec:cortex-learning}
% --------------------------------------------------------------------------

The self-ingestion loop of \S\ref{sec:cortex-self} reveals the system's
output fidelity at a given moment, but does not by itself constitute learning.
To learn, the system must \textbf{change which HLLSets compose its current
state} in response to accumulated evidence. The IICA constraint---HLLSets are
immutable once created, and $\text{BSS}_{\tau}$, $\text{BSS}_{\rho}$
between fixed HLLSets are computed values that cannot be modified---forces a
specific kind of learning: \textbf{structural selection through threshold
tuning.}

\subsubsection{The Two Thresholds}

Let $\mathcal{S}(t) = \{H_1, H_2, \ldots, H_k\}$ be the collection of
HLLSets representing the system's current state at time $t$---the working set
of page-level, chapter-level, and self-ingested HLLSets that feed into the
materializer and decoder. Learning operates on $\mathcal{S}$ through two
thresholds:

\begin{definition}[Inclusion Threshold $\theta_{\tau}$]
An HLLSet $H_i$ is admitted to the current state collection only if its
forward inclusion with the source document exceeds the threshold:
\begin{equation}
  H_i \in \mathcal{S}(t) \iff \text{BSS}_{\tau}(H_i \to H_{\text{doc}})
  \geq \theta_{\tau}.
\end{equation}
$\theta_{\tau}$ filters out low-fidelity representations---measurements
that lost too much source structure during encoding, gating, or
materialization.
\end{definition}

\begin{definition}[Divergence Threshold $\theta_{\rho}$]
An HLLSet $H_i$ is excluded from the current state collection if its reverse
exclusion with the source document exceeds the threshold:
\begin{equation}
  H_i \notin \mathcal{S}(t) \iff \text{BSS}_{\rho}(H_i \to
  H_{\text{doc}}) > \theta_{\rho}.
\end{equation}
$\theta_{\rho}$ filters out divergent representations---measurements
containing structure absent from the source, i.e., encoding noise, gate
leakage, or materializer hallucination.
\end{definition}

The two thresholds together define a \textbf{validity region} in
BSS-space. HLLSets within the region $(\tau \geq \theta_{\tau}) \land
(\rho \leq \theta_{\rho})$ compose the system's operative world model;
HLLSets outside it are excluded from driving output but remain
content-addressed and addressable.

\subsubsection{Threshold Tuning as the Cortical Brake in Operation}

The learning process adjusts $\theta_{\tau}$ and $\theta_{\rho}$ over
time. This is not parametric optimization (there is no loss function, no
gradient, no weight update). It is \textbf{structural regulation}: the
Noether controller monitors the cross-layer Fisher matrix
$F_{bb'}$ across the temporal pyramid and detects when the current threshold
settings produce drift, oscillation, or collapse.

Three regimes are distinguishable:

\begin{enumerate}[leftmargin=*,nosep]
  \item \textbf{Undersuppression} ($\theta_{\tau}$ too low, $\theta_{\rho}$
    too high). The state collection $\mathcal{S}$ admits low-fidelity and
    divergent HLLSets. The decoder output oscillates---successive
    materializations of the same document produce different token sequences.
    In Lincoln's terms, this is the cortical brake failing to engage: the
    substrate-level reinforcement signal (TF accumulation) drives behavior
    without adequate regulatory filtering.
  \item \textbf{Oversuppression} ($\theta_{\tau}$ too high, $\theta_{\rho}$
    too low). $\mathcal{S}$ is nearly empty. The system has excluded almost
    all HLLSets from its working set. Output is sparse or null. In Lincoln's
    terms, this is pathological over-engagement of the habenular pathway---
    the system suppresses action even when evidence supports it, analogous to
    the anhedonia and motivational deficits associated with LHb hyperactivity
    in depression \cite{lincoln2026}.
  \item \textbf{Regulated operation} ($\theta_{\tau}$ and $\theta_{\rho}$
    at values that maintain $\mathcal{S}$ at a stable size with low output
    variance across temporal layers). The cross-layer Fisher matrix shows
    convergence: the same HLLSets persist in $\mathcal{S}$ across successive
    temporal windows, and the BSS between successive outputs is high and
    stable. This is the EWM analog of adaptive epistemic functioning under
    intact cortical regulation.
\end{enumerate}

\subsubsection{The Matching Law Revisited: Attention as Inclusion}

The threshold mechanism gives operational meaning to the matching-law analysis
of \S\ref{sec:intrinsic-reward}. The system's ``attention allocation''
across HLLSets is not a continuous weight vector but a \textbf{binary
inclusion decision}: each $H_i$ is either in $\mathcal{S}$ (it drives output)
or out (it does not). The proportion of HLLSets from a given measurement
strategy (e.g., a specific bootstrap family, a specific gate, a specific
temporal window) that survive the thresholds is the EWM analog of the matching
law's $B_1 / B_2$ ratio.

The sensitivity parameter $s$ maps onto the \textbf{steepness of the
threshold boundary}: a sharp boundary ($s$ high) means small differences in
BSS produce large differences in inclusion probability; a soft boundary ($s$
low) means the system is tolerant of variance. The bias parameter $b$ maps
onto the \textbf{relative positioning} of $\theta_{\tau}$ and
$\theta_{\rho}$: a system biased toward inclusivity sets $\theta_{\tau}$
low and $\theta_{\rho}$ high; a system biased toward conservatism does the
opposite.

\subsubsection{Exclusion Without Erasure}

The critical structural consequence of learning-as-threshold-tuning is that
\textbf{excluded HLLSets are not deleted}. They are immutable and
content-addressed; their CIDs remain valid; their TF contributions to the LUT
persist. The system ``forgets'' nothing---it only \textbf{stops attending.}

This has two implications:

\begin{enumerate}[leftmargin=*,nosep]
  \item \textbf{Reversible exclusion.} Lowering $\theta_{\tau}$ (or raising
    $\theta_{\rho}$) re-admits previously excluded HLLSets to $\mathcal{S}$
    instantly, at their earned TF rank. No recomputation, no re-measurement, no
    cold start. The evidence was always stored; only the regulatory threshold
    changed. This is the structural realization of Lincoln's account of belief
    revision: when the cortical framework changes, evidence that was always
    present but previously inadmissible becomes salient---not because the
    evidence changed, but because the threshold for admission changed.
  \item \textbf{No catastrophic forgetting.} Because exclusion is reversible,
    the system cannot permanently lose knowledge through threshold drift. A
    pathological oscillation that excludes and re-admits HLLSets in rapid
    succession produces unstable output but does not corrupt the underlying
    lattice. The IICA property guarantees that the same $H_i$ is the same
    $H_i$ regardless of how many times it cycles in and out of $\mathcal{S}$.
\end{enumerate}

This stands in marked contrast to neural network learning, where weight
updates are destructive (the old weight is overwritten) and catastrophic
forgetting is a persistent problem. In the EWM framework, learning is
\textbf{attention, not erasure}---the structural analog of Lincoln's
observation that the NAc-LHb reinforcement signal itself never changes; only
the cortical regulation of which signals reach behavioral output changes.

\end{document}
